\chapter{Representations of the time harmonic elastic wave Green's 
tensor}\label{Grep}
%\ \hrulefill
%\rule{154.0mm}{10pt}\\
We recall from chapter \ref{MSrep} that
 the material 
properties of the propagation environment have been defined as  
the sum of a reference and residual component, i.e.
$(\rho, \alpha, \beta)=(\rf{\rho}+{\mit \Delta}\rho,\rf{\alpha}+{\mit 
\Delta}\alpha,\rf{\beta}+{\mit \Delta}\beta)$. In this chapter, we will 
be concerned exclusively with uniform media so that 
residual properties $({\mit \Delta}\rho,{\mit 
\Delta}\alpha,{\mit \Delta}\beta)$ vanish and we have
$(\rho, \alpha, \beta)=(\rf{\rho},\rf{\alpha},\rf{\beta})$. To simplify 
notation, we discontinue the practice of writing  $(\rf{\rho},\rf{\alpha},\rf{\beta})$
to denote reference medium parameters and write 
$({\rho},{\alpha},{\beta})$ instead.

This chapter is concerned with the construction of the time harmonic 
elastic wave motion driven by a {\em lattice of multipolar point forces\/} 
${\bmi f}$ {\em  superimposed on a homogeneous medium\/}.
The homogeneous medium is
characterized by the  elastic parameters $(\rho,\alpha,\beta)$, which 
denote the density, compressional and shear speeds respectively, in an 
isotropic medium. 

Each point force  ${\bmi f}$,  has an assumed time 
dependence of the form \linebreak
$\exp(\epsilon t -i 2 \pi \nu t)$ where 
$\nu $ is the  frequency and $\epsilon$ is a non-negative {\em 
growth\/} parameter. The {\em adiabatic growth term\/} ${\rm e}^{\epsilon 
t}$, 
evolves the wavefield from a state of quiescence in the infinite past, to 
its 
present value, following a path that ensures that the {\em wavefield is 
always in steady 
state equilibrium with the driving force\/} ${\bmi f}$. This device has been 
employed by 
{\em Lighthill} ~\shortcite{Lighthill} to enforce the {\em radiation 
boundary conditions\/}.
For  convenience, we introduce the {\em complex angular frequency\/}
\begin{equation} \label{w+}
\omega\equiv 2 \pi \nu +i \epsilon.
\end{equation}
We can then synthesize {\em transient motion\/} from a discrete set of 
{\em 
forced motions\/}, sampled from the {\em Bromwich contour\/} in the 
complex frequency 
domain [e.g.{\em Chapman} ~\shortcite{Chap78}] as
\begin{equation}	\label{InvLaplace}
	\left[{\cdot }\right](t)=2 \, {\rm e}^{\epsilon  t}\,
	\Re \int_{ 0}^{\nu_{max}} {\rm d}\nu \ {\rm e}^{ -i 
	2 \pi \nu t}\left[{\cdot }\right] (2 \pi \nu +i \epsilon)
\end{equation}
where $\nu_{{max}}$ is the maximum frequency of the source field. 
Usually the growth parameter $\epsilon$ is chosen so that $5 \le \epsilon T
\le 10$ \cite{Banerjee}, where $T$ is the length of the required time series (which is 
usually determined by trial and error). Using the complex angular frequency $\omega$ smoothes the frequency domain 
response, which enhances the conditioning of the problem, 
causing the complex frequency multiple scattering series to converge 
faster than the real frequency version. However, we must be aware that the 
process is equivalent to an analytic continuation, and is inherently 
unstable. The instability, is manifested by the sensitivity of 
(\ref{InvLaplace}) to errors in the frequency domain response. This is a 
particularly troubling issue, as the multiple scattering series is 
usually only summed to a relative error of $10^{-6}$ and it remains an 
open question whether we have secured sufficient accuracy to avoid the 
deleterious effects of the numerical instability inherent in 
(\ref{InvLaplace}). An excellent survey of numerical methods for the 
numerical evaluation of (\ref{InvLaplace}) is presented in 
{\em Davies \& Martin\/}~(1979).



%%%%%%%%%%%%%%%%%%%%%%%%%%%%%%%%%%%%%%%%%%%%%%%%%%%%%%%%%%%%%%%%%%%%%%%%%%%%%%%%%%


\section{The free space, time harmonic,  elastic wave Green's tensor} \label{classpwprop}
Given the 2-D vectors
\[
{\bmi  x}\equiv (z,x), \quad {\bmi  x}^{\prime}\equiv (z^{\prime},x^{\prime})
\]
a particularly useful solution $\rf{\bmi  G }({\bmi  x};{\bmi  x}^{\prime};k_y)$, 
satisfies (\ref{govEqn0}) 	when the the material properties in 
Table \ref{ATable} are homogeneous and the  {\em $y-$Fourier transformed body 
force\/} ${\bmi F}({\bmi x};{\bmi x}^{\prime})$ assumes the 
special form 
\[
{\bmi F}({\bmi x};{\bmi x}^{\prime})=
{\bmi  I}_6 \delta({\bmi x}-{\bmi x}^{\prime}),
\] 
with ${\bmi  I}_6$ the $6 \times 6$ {\em identity matrix\/}, {\em viz\/}
\begin{equation} \label{(2.1.6)} 
{{\CB L}}\cdot \rf{\bmi  G}({\bmi  x};{\bmi  x}^{\prime};k_y)={\bmi  I}_6 
\delta({z}-{z}^{\prime})\delta({x}-{x}^{\prime}).
\end{equation}
We can then use the linearity of (\ref{govEqn0}) to pose the {\em 
particular solution\/} ${\bmi  b}_{P}$ of (\ref{govEqn0}) as the {\em 
superposition integral\/}
\begin{equation}	\label{superposition}
	{\bmi  b}_{P}=\int {\rm d}{\bmi  x}^{\prime}\, \rf{\bmi  G}({\bmi  x};{\bmi 
x}^{\prime};k_y){\bmi 
	f}({\bmi  x}^{\prime})
\end{equation}
for an arbitrary body force $\bmi  f$.
It will be  convenient for the purpose of subsequent constructions, to 
interpret 
$\rf{\bmi  G}({\bmi  x};{\bmi  x}^{\prime};k_y)$  of (\ref{superposition}), as the 
{\em integral kernel\/} of the formal inverse
\begin{equation}	\label{invDef}
	\left[{{\CB L}}^{-1} {\bmi  w}\right]({\bmi  x})\equiv \int {\rm d}{\bmi 
	x}^{\prime}\, \rf{\bmi  G }({\bmi  x};{\bmi  x}^{\prime};k_y)\cdot{\bmi  w}({\bmi 
x}^{\prime}) 
\end{equation}
for an arbitrary {\em test\/} vector-function $\bmi  w$, where ${{\CB L}}$ is the governing 
operator in (\ref{govEqn0}). Equation 
(\ref{invDef}) is strictly formal,  as the {\em null space\/} of 
${{\CB L}}$ is spanned by the set of {\em homogeneous\/} solutions and 
consequently, 
$\rf{\bmi  G}$ is not unique. To get a {\em well posed\/} inverse, we invoke 
the {\em radiation condition\/} that {\em there are no homogeneous 
solutions 
of\/} (\ref{govEqn0}) {\em due to sources at infinity\/}.



To solve for $\rf{\bmi  G}$ in (\ref{(2.1.6)}), or equivalently, to construct 
the 
inverse  ${{\CB L}}^{-1}$ in (\ref{invDef}), we follow the 
classical {\em Fourier Transform\/} route. To begin, we transform the $x-$ 
dependence via ${\cal F}_{x}\, [\cdot]({k}_{x})\equiv \frac{1}{\sqrt{2 
\pi}}\int_{-\infty }^{ \infty 
}dx\ \exp \left({-i\, {k}_{x} x}\right)\, [\cdot](x)$, where  ${k}_{x}$ is the 
horizontal 
component 
of the wavevector,
so that the transverse {\em Fourier\/} transform of the governing equation 
(\ref{(2.1.6)})
may be written as
\begin{eqnarray}	\label{FLdef}
\lefteqn{\langle { {k}_{x}} | {\CB L }| { {k}_{x}}^{\prime}\rangle
\langle { {k}_{x}}^{\prime}| x\rangle\langle x| \,\rf{\bmi 
G}({\bmi x};{\bmi x}^{\prime};k_y)}\nonumber\\
&&\equiv{\bmi  L}(z,{k}_{x},{k}_{y}) \,\rf{\bmi  G}(z,{k}_{x};z^{\prime},x^{\prime};k_y)\nonumber\\
&& \equiv \left[\partial_z -{\bmi  A}({k}_{x},{k}_{y})\,  \right]\,\rf{\bmi 
G}(z,{k}_{x};z^{\prime},x^{\prime};k_y)={\bmi  I}_{6} 
	\delta(z-z^{\prime}) \, \frac{\exp(-i {k}_{x} x^{\prime})}{\sqrt{2 \pi}}
\end{eqnarray}
where the {\em Fourier domain system matrix\/} $\langle { {k}_{x}} | {\CB A }| 
{ {k}_{x}}^{\prime}\rangle= \delta({ {k}_{x}}-{ {k}_{x}}^{\prime}) {\bmi  
A}({k}_{x},{k}_{y})$ is defined
\begin{equation}	\label{Akdef}
	{\bmi  A}({k}_{x},{k}_{y})=\left[{\begin{array}{cccccc} 
0&-i \gamma {  k}_{  x}& 
 - \gamma   i\ {k}_{y}&a&0&0\\ -i{k}_{x}&0&0&0&b&0\\ 
-i{k}_{y}&0&0&0&0&b\\ - \rho {\omega }^{  2}&  
0&0&0&-i{k}_{x}&-i{k}_{y}\\ 
0&- \rho {\omega }^{  2}  + 
\zeta {  k}_{  x}^{  2}  + \mu {  
k}_{  y}^{  2}&{  k}_{  x}{  k}_{ 
 y}\left({\chi   + \mu }\right)&  -i{k}_{x} \gamma 
&  0&0\\ 0&{k}_{x}{k}_{y}\left({ \chi   + \mu 
}\right)&- \rho {\omega }^{  2}  + \zeta {  
k}_{  y}^{  2}  + \mu {  k}_{  
x}^{  2}&  -i{k}_{y} \gamma &  0&0\end{array}}\right]
\end{equation}
due to the {\em translational invariance of\/} ${\CB A }$.
The system operator ${\bmi  A}({k}_{x},{k}_{y})$ of 
(\ref{Akdef}) is {\em quasi-Hamiltonian\/} since
\begin{equation} \label{ATdef}
{\bmi  A}^{T}(-{k}_{x},-{k}_{y}){\bmi  J}_6 + {\bmi  J}_6 {\bmi  
A}({k}_{x},{k}_{y}) \equiv {\bf 0}_6
\end{equation}
where ${\bmi  J}_6$ is the {\em canonically symplectic\/} array of 
(\ref{Jdef}).
The {\em diagonal array\/} ${\bmi 
\Lambda}$ {\em of eigenvalues\/}  of ${\bmi  A}({k}_{x},{k}_{y})$ is defined as
\begin{samepage}\begin{eqnarray}	\label{eigDef}
\lefteqn{	{\bmi \Lambda}({k}_{x},{k}_{y}) =}\\
&&\!\!\!\!\!\!\!\!\!\!\!\!
i \ \mbox{diag} \left[{k}_{z,P 
	}({k}_{x},{k}_{y}),{k}_{z,S }({k}_{x},{k}_{y}),{k}_{z,H }({k}_{x},{k}_{y}),
	-{k}_{z,P }({k}_{x},{k}_{y}),-{k}_{z,S }({k}_{x},{k}_{y}),
	-{k}_{z,H }({k}_{x},{k}_{y}) \right] \nonumber
\end{eqnarray}\end{samepage}
where
\begin{equation}	\label{kzcDef}
	{k}_{z,c}({k}_{x},{k}_{y}) = \left\{
	\begin{array}{c}
	\sqrt {\hat{ K}_{c}^{2}({k}_{y})\ -\ {k}_{x}^{2}}\ \ \ \ \ \ |{k}_{x}|\ <\ 
	|\hat{K}_{c}({k}_{y})|\\
	i\sqrt {{k}_{x}^{ 2}\ -\ \hat{ K}_{c}^{2}({k}_{y})} \ \ \ \ |{k}_{x}|\ \ge   \ 
	|\hat{K}_{c}({k}_{y})|
\end{array}
\right. , \quad  c\in\{P,S,H\}, 
\end{equation}
with
\begin{equation}	\label{Kdef}
	\hat{K}_{c}({k}_{y}) = \left\{
	\begin{array}{c}
	\sqrt {(\omega/{c})^{2}\ -\ {k}_{y}^{2}}\ \ \ \ \ \ |{k}_{y}|\ <\ 
	|\omega/{c}|\\
	i\sqrt {{k}_{y}^{ 2}\ -\ (\omega/{c})^{2}} \ \ \ \ |{k}_{y}|\ \ge   \ |\omega/{c}|
\end{array} \right. , \quad  c\in\{P,S,H\}. 
\end{equation}
{\em We note that ${k}_{z,S }={k}_{z,H }$ 
because $S=H=\beta$ from Table \ref{ATable}. We distinguish 
between ${k}_{z,S }$ and ${k}_{z,H }$,  only to maintain a self-consistent 
notational convention.\/}
 The system matrix ${\bmi  A}({k}_{x},{k}_{y})$ of (\ref{Akdef}) has an equivalent {\em 
spectral representation\/}:
\begin{equation}	\label{specA}
	{\bmi  A}({k}_{x},{k}_{y}) \equiv {\bmi  D}_{z}({k}_{x},{k}_{y}){\bmi 
	\Lambda }({k}_{x},{k}_{y}){\bmi  D}_{z}^{ -1}({k}_{x},{k}_{y});
\end{equation} 
where the {\em eigen matrix\/} ${\bmi  D}_{z}({k}_{x},{k}_{y})$ is  defined as
\begin{eqnarray} \label{Ddef}
\lefteqn{{\bmi D}_{z}(k_{x},k_{y})=}\\
&&\!\!\!\!\!\!\!\!\!\!\!\!
\left[\begin{array}{cccccc}
\left|+,P, k_{x},k_{y}  \right\rangle&\left|+,S, k_{x},k_{y}  
\right\rangle&\left|+,H, k_{x},k_{y}  \right\rangle&
\left|-,P, k_{x},k_{y}  \right\rangle&\left|-,S, k_{x},k_{y}  
\right\rangle&\left|-,H, k_{x},k_{y}  \right\rangle
\end{array}\right]\nonumber
\end{eqnarray}
with:
\begin{itemize}
\item the $\left|\pm,P, k_{x},k_{y}  \right\rangle$ eigen-vectors 
are the   ${\stackrel{\mbox{\small downgoing}}{{}_{\mbox{\small 
upgoing}}}}$  $P-$waves 
\begin{equation} \label{Peigen}
\left|\pm,P, k_{x},k_{y}  \right\rangle=\left[\begin{array}{c} \pm 
i\,{k}_{z,P}(k_x,k_y)\\i\,k_{x}\\i\,k_{y}\\
  \rho \,\left( 2\,{\beta^2}\,{k_{x}^2} + 2\,{\beta^2}\,{k_{y}^2} - 
     {\omega^2} \right) \\ \mp 2\,\rho\,{\beta^2}\,k_{x}\,{k}_{z,P}(k_x,k_y) \\
  \mp 2\,\rho\,{\beta^2}\,k_{y}\,{k}_{z,P}(k_x,k_y) \end{array}\right]
  \epsilon_{P}
\end{equation}
with the  normalization factor ${\epsilon}_{P}$

\item the $\left|\pm,S, k_{x},k_{y}  \right\rangle$ eigen-vectors 
are the   ${\stackrel{\mbox{\small downgoing}}{{}_{\mbox{\small 
upgoing}}}}$  quasi-$SV-$waves 
\begin{equation} \label{SVeigen}
\left|\pm,S, k_{x},k_{y}  \right\rangle=\left[\begin{array}{c} 
i\,k_{x}\\
\mp i\,{k}_{z,S}(k_x,k_y)\\0\\
\mp 2\,\rho\,{\beta^2}\,k_{x}\,{k}_{z,S}(k_x,k_y) \\
  \rho \,\left({\omega^2} -2\,{\beta^2}\,{k_{x}^2} - {\beta^2}\,{k_{y}^2}  \right) \\
  -\rho\,{\beta^2}\,k_{x}\,k_{y} \end{array}\right]
  \epsilon_{S}
\end{equation}
with the  normalization factor ${\epsilon}_{S}$

\item the $\left|\pm,H, k_{x},k_{y}  \right\rangle$ eigen-vectors 
are the   ${\stackrel{\mbox{\small downgoing}}{{}_{\mbox{\small 
upgoing}}}}$  quasi-$SH-$waves 
\begin{equation} \label{SHeigen}
\left|\pm,H, k_{x},k_{y}  \right\rangle=\left[\begin{array}{c} 
\mp k_{y}\,{k}_{z,H}(k_x,k_y) \\-k_{x}\,k_{y} \\
  \hat{K}_{H}^2\\
2\,i\,k_{y}\,\rho \,\left({\beta^2}\,{k_{x}^2} + 
     {\beta^2}\,{k_{y}^2} - {\omega^2} \right) \\
  \mp 2\,i\,\rho\,{\beta^2}\,k_{x}\,k_{y}\,{k}_{z,H}(k_x,k_y) \\
  \pm i\,{k}_{z,H}(k_x,k_y)\,\rho \,\left(  {\omega^2} -2\,{\beta^2}\,{k_{y}^2} \right) \end{array}\right]
  \epsilon_{H}
\end{equation}
with the  normalization factor ${\epsilon}_{H}$
\end{itemize}
We have called the eigen-vectors in (\ref{SVeigen}) and (\ref{SHeigen}) 
{\em quasi-\/}$SV$ and {\em quasi-\/}$SH$ waves respectively, because the 
corresponding vector fields are {\em solenoidal\/} and are {\em in the 
$(z,x)$ plane\/} 
and {\em out of the $(z,x)$ plane\/} respectively, when $k_{y}=0$. However, our 
choice of the  shear wavefield decomposition into quasi-inplane and quasi-out 
of plane components
differs from the choice commonly used in seismology [e.g. 
{\em Prange \& Toks\"oz}~(1990)]. Unlike the conventional decomposition,  the basis vectors 
(\ref{SVeigen}) and (\ref{SHeigen}) are {\em non-degenerate\/} when 
${k}_{x}=0$ and ${k}_{y}=0$.
The normalization factors ${\epsilon}_{P},{\epsilon}_{S} $ and 
${\epsilon}_{H}$
in Eqns.(\ref{Peigen})-(\ref{SHeigen}) are chosen so that ${\bmi  
D}_{z}^{-1}({k}_{x},{k}_{y})$ 
takes a simple form.

%%%%%%%%%%%%%%%%%%%%%%%%%%%%%%%%%%%%%%%%%%%%%%%%%%%%%%%%%%%%%%%%%%%%%%%%%%%%%%%%%%%%%%%%%
%%%%%%%%%%%%%%%%%%%%%%%%%%%%%%%%%%%%%%%%%%%%%%%%%%%%%%%%%%%%%%%%%%%%%%%%%%%%%%%%%%%%%%%%%
%%%%%%%%%%%%%%%%%%%%%%%%%%%%%%%%%%%%%%%%%%%%%%%%%%%%%%%%%%%%%%%%%%%%%%%%%%%%%%%%%%%%%%%%%
%%%%%%%%%%%%%%%%%%%%%%%%%%%%%%%%%%%%%%%%%%%%%%%%%%%%%%%%%%%%%%%%%%%%%%%%%%%%%%%%%%%%%%%%%
%%%%%%%%%%%%%%%%%%%%%%%%%%%%%%%%%%%%%%%%%%%%%%%%%%%%%%%%%%%%%%%%%%%%%%%%%%%%%%%%%%%%%%%%

Using the standard result that eigenspaces of ${\bmi  A}^{T}({k}_{x},{k}_{y})$ annihilate
eigenspaces of ${\bmi  A}({k}_{x},{k}_{y})$ (and vice versa) 
[e.g. {\em Milne\/}~\shortcite{Milne}, 
{\em p.89\/}], we can avoid constructing  ${\bmi  D}_{z}^{-1}({k}_{x},{k}_{y})$ in (\ref{specA}) by explicit 
matrix inversion, since we use the {\em Hamiltonian symmetry\/} 
(\ref{ATdef}), to relate ${\bmi  D}_{z}^{-1}({k}_{x},{k}_{y})$ to the eigen matrix of the 
{\em transposed\/} or {\em 
dual\/} system matrix ${\bmi  A}^{T}({k}_{x},{k}_{y})$ defined as
\[
{\bmi  A}^{T}({k}_{x},{k}_{y})\equiv {\bmi  J}_6 {\bmi  A}(-{k}_{x},-{k}_{y}) {\bmi  J}_6
\] 
To begin, it follows from (\ref{eigDef}) and (\ref{kzcDef}) that ${\bmi 
\Lambda}(-{k}_{x},-{k}_{y})={\bmi \Lambda} ({k}_{x},{k}_{y})$, 
so that
\begin{samepage}
\begin{eqnarray*}
{\bmi  D}_{z}(-{k}_{x},-{k}_{y}){\bmi \Lambda} ({k}_{x},{k}_{y})&=&
{\bmi  A}(-{k}_{x},-{k}_{y}){\bmi  D}_{z}(-{k}_{x},-{k}_{y})\\
&=&{\bmi  J}_{6}{\bmi 
A}^{T}({k}_{x},{k}_{y}){\bmi  J}_{6}{\bmi  D}_{z}(-{k}_{x},-{k}_{y})
\end{eqnarray*}
\end{samepage}
which implies upon rearrangement that
\[
{\bmi  A}^{T}({k}_{x},{k}_{y}){\bmi  J}_{6}{\bmi  D}_{z}(-{k}_{x},-{k}_{y})=-{\bmi  J}_{6}
{\bmi  D}_{z}(-{k}_{x},-{k}_{y}){\bmi  
\Lambda} ({k}_{x},{k}_{y}).
\]
Consequently, ${\bmi  J}_{6}{\bmi  D}_{z}(-{k}_{x},-{k}_{y})$ is an eigenbasis of the dual 
operator ${\bmi 
A}^{T}({k}_{x},{k}_{y})$, with $-{\bmi  \Lambda} ({k}_{x},{k}_{y})$ the corresponding eigenvalue array. 
From direct evaluation it follows that
\begin{equation} \label{dinv1}
{\left({{\bmi  J}_{6}{\bmi  D}_{z}(-{k}_{x},-{k}_{y})}\right)}^{T}{\bmi  
D}_{z}({k}_{x},{k}_{y})= i\left[{\begin{array}{cc}{\bf 0}_{3}&{\bmi V}\\
-{\bmi V}&{\bf 0}_{3}\end{array}}\right]
\end{equation}
where
\[
{\bmi V}=
\mbox{diag}\left[\begin{array}{ccc}
2 \rho {\omega }^{  2}{  k}_{  z, P }({k}_{x}){\epsilon 
}_{P}^{  2},\,&
2 \mu {  k}_{  z, S }({k}_{x}) \hat{K}_{S }^{2}{\epsilon }_{S}^{  
2},\,&
2 \rho {\omega }^{  2}{  k}_{  z, H }({k}_{x}) \hat{K}_{H }^{2}{\epsilon }_{H}^{  2}
\end{array}\right]
\]
and  $\hat{K}_{c}, c\in\{P,S,H\}$ was given by (\ref{Kdef}),
 so that if we choose the normalization factors ${\epsilon }_{P }$, 
${\epsilon }_{S}$ and ${\epsilon }_{H}$ in Eqns.(\ref{Peigen})-(\ref{SHeigen})
according to 
 \begin{equation} \label{epsdef}
{\epsilon }_{P }= \frac{1}{\sqrt{2 \rho {\omega }^{  2}{  k}_{  
z, P }}}, \quad
{\epsilon }_{S }= \frac{\omega}{\beta\,\hat{K}_{S }\,\sqrt{2 \rho {\omega }^{  2}  {  k}_{  
z, S }}}
, \quad \mbox{and} \quad
{\epsilon }_{H }= \frac{1}{\hat{K}_{H }\, \sqrt{2 \rho {\omega }^{  2}  {  k}_{  
z, H }}} 
\end{equation}
 then (\ref{dinv1}) simplifies to
\begin{equation} \label{dinv2}
{\left({{\bmi  J}_{6}{\bmi  D}_{z}(-{k}_{x},-{k}_{y})}\right)}^{T}{\bmi  D}_{z}({k}_{x},{k}_{y})
=i {\bmi  J}_{6}
\end{equation}
so that upon rearrangement of (\ref{dinv2}) we get
\begin{equation}	\label{D1def}
{\bmi  D}_{z}^{-1}({k}_{x},{k}_{y})=-i\	{\bmi  J}_6{\bmi  D}_{z}^{T}(-{k}_{x},-{k}_{y}){\bmi  J}_6
\end{equation}
Equation (\ref{D1def}) for the inverse is reminiscent of the {\em 
symplectic\/} structures found in {\em Hamiltonian dynamical systems\/} 
[e.g. {\em 
Goldstein}~\shortcite{Goldstein}] ({\em see also \/} {\sc Box} 
$\bf 3.1$). 



%%%%%%%%%%%%%%%%%%%%%%%%%%%%%%%%%%%%%%%%%%%%%%%%%%%%%%%%%%%%%%%%%%%%%%%%%%%%%%%%%%%%%
%%%%%%%%%%%%%%%%%%%%%%%%%%%%%%%%%%%%%%%%%%%%%%%%%%%%%%%%%%%%%%%%%%%%%%%%%%%%%%%%%%%%%
\begin{figure}
\fbox{
\begin{minipage} {142.0mm}
{\sc Box} $\bf 3.1$ \newline
A {\em fundamental\/} investigation of 
the seismic wave equations as a {\em Hamiltonian dynamical system\/} may 
be found 
in {\em Kennett }~\shortcite{K74b},
{\em Chapman \& Woodhouse }~\shortcite{ChapWood}, where it is shown that 
the coupled 
second order elastic wave equations are the {\em Euler\/} equations for 
the stationary value of a {\em Lagrangian variational principle\/}. 
The ``Hamiltonian" first order system employed in this thesis is then 
obtained via a {\em Legendre transformation\/} [e.g. {\em 
Goldstein}~\shortcite{Goldstein} Ch. $12$]. Unlike, the equations of 
{\em static elastic equilibrium\/}, variational principles for the time 
harmonic elastic wave equations are not {\em extremal\/}, due to the 
complex nature of the {\em radiation conditions\/}, which siphon energy 
out of the system, as well as the {\em D'Alembert forces\/}, which render 
the wave equations {\em indefinite\/}. The lack of an extremum, vitiates 
some of the numerical procedures that have been developed by 
mathematicians and engineers to handle the static problem. None the less,
identifying the time harmonic wave elastic wave equation as a Hamiltonian 
dynamical 
system suggests a very useful {\em propagation invariant\/} that follows 
from the {\em incompressibility of a Hamiltonian flow in  
displacement-traction phase space\/} [e.g. {\em 
Schultz\/}~\shortcite{Schultz}], so that 
${\bmi  b}_{ 1}\wedge {\bmi  b}_{ 2}|_{z_1}={\bmi  b}_{ 1}\wedge {\bmi  b}_{ 
2}|_{z_2}$. In theoretical seismology, the exterior product ``$\wedge$" is 
expressed in terms of the canonically symplectic array ${\bmi  J}_6$ of (\ref{Jdef}), 
so that
\[
{\bmi  b}_{ 1}^{ T} (z_1,-{k}_{x},-{k}_{y}){\bmi  J}_{6}{\bmi  b}_{2}(z_1,{k}_{x},{k}_{y})=
{\bmi  b}_{ 1}^{ T} (z_2,-{k}_{x},-{k}_{y}){\bmi  J}_{6}{\bmi  
b}_{2}(z_2,{k}_{x},{k}_{y})\hspace{9 em}(1)
\]
for a pair of wavefields ${\bmi  b}_{1},{\bmi  b}_{2}$ in a {\em source 
free\/} depth interval $z_1 \le \zeta \le z_2$.
An appreciable portion of the {\em matrix  mechanics\/} for seismic wave 
propagation in stratified 
media, may be developed, with (1) as the starting point 
\cite{K83}. The identity (1) plays a role analogous to the 
{\em Rayleigh-Betti\/} representation of the conventional coupled 
second order theory. 
\vspace{6mm}
\end{minipage}
}
\end{figure}
%%%%%%%%%%%%%%%%%%%%%%%%%%%%%%%%%%%%%%%%%%%%%%%%%%%%%%%%%%%%%%%%%%%%%%%%%%%%%%%%%%%%%

%%%%%%%%%%%%%%%%%%%%%%%%%%%%%%%%%%%%%%%%%%%%%%%%%%%%%%%%%%%%%%%%%%%%%%%%%%%%%%%%%%%%%



We now transform the $z-$dependence of  (\ref{FLdef}) using
the {\em  Fourier transform\/} 
\[
{\cal F}_{z}\, 
[\cdot] ({k}_{z}) \equiv \frac{1}{\sqrt{2 \pi}}\int_{-\infty }^{ \infty } 
dz\ \exp \left({-i\ {k}_{z}\ z}\right)\, [\cdot] (z) ,\]
 to get
\begin{equation}	\label{FzLdef}
	\tilde{\bmi  L}({\bmi k}) \,\rf{\bmi  G}({\bmi k};{\bmi x}^{\prime};k_y) \equiv \left[i 
{\bmi  I}_6 {k}_{z} -{\bmi 
A}({k}_{x},{k}_{y})\,  
	\right]\,\rf{\bmi  G}({\bmi k};{\bmi x}^{\prime};k_y)={\bmi  I}_{6} 	
	\frac{\exp({-i {\bmi k}\cdot {\bmi x}^{\prime}})}{{2 \pi}},
\end{equation}
where ${\bmi x}^{\prime}=(z^{\prime},x^{\prime})$ and  ${\bmi 
k}=(k_{z},k_{x})$, so that
\begin{equation}	\label{FzLdef1}
	\rf{\bmi  G}({\bmi k};{\bmi x}^{\prime};k_y) \equiv
	\left[i {\bmi  I}_6 {k}_{z} -{\bmi  A}({k}_{x},{k}_{y}) \right]^{-1} 
	\frac{\exp({-i {\bmi k}\cdot {\bmi x}^{\prime}})}{{2 \pi}}.
\end{equation}
Using the spectral representation (\ref{specA}) in  (\ref{FzLdef1}), 
we get 
\begin{eqnarray} \label{InvLDef0}
\rf{\bmi  G}({\bmi k};{\bmi x}^{\prime};k_y) 
&=&{\bmi  D}_{z}({k}_{x},{k}_{y}) 
{\left[i\, {\bmi  I}_{6}\, {k}_{z}\, -\, i\,   {\bmi \Lambda} \right]}^{-1}
{\bmi  D}_{z}^{-1}({k}_{x},{k}_{y}) \frac{\exp({-i {\bmi k}\cdot {\bmi x}^{\prime}})}{2 \pi}
\nonumber\\
&=&{\bmi  D}_{z}({k}_{x},{k}_{y}) 
{\left[{\bmi  J}_{6}\, {k}_{z}\, - \,{\bmi  J}_{6}\, {\bmi \Lambda} 
\right]}^{-1}
{\bmi  D}_{z}^{T}(-{k}_{x},-{k}_{y})\,{\bmi  J}_6
\frac{\exp(-i {\bmi k}\cdot {\bmi x}^{\prime})}{2 \pi}\nonumber\\
&=&{\bmi  D}_{z}({k}_{x},{k}_{y})
\left[{\begin{array}{cc}{\bf 0}_{3}&-{\bmi W}^{+}\\
{\bmi W}^{-}&{\bf 0}_{3}\end{array}}\right]
{\bmi  
D}_{z}^{T}(-{k}_{x},-{k}_{y})
\,{\bmi  J}_6 \frac{\exp(-i {\bmi k}\cdot {\bmi x}^{\prime})}{{2 \pi}}
\end{eqnarray} 
where
\[
{\bmi W}^{\pm}=\mbox{diag}
\left[\begin{array}{ccc}
\frac{1}{({k}_{z}\mp{k}_{z,P })},\,&\frac{1}{({k}_{z}\mp{k}_{z,S 
})},\,&\frac{1}{({k}_{z}\mp{k}_{z,S })}\end{array}\right]
\]


Applying the inverse Fourier transform 
${\CB F}^{-1}\, [\cdot]({\bmi x}) \equiv \frac{1}{{2 \pi}} 
\int_{R^2 } {\rm d}{\bmi k} \exp \left(i {\bmi k}\cdot {\bmi x}\right)\, 
[\cdot] ({\bmi k})$, 
it follows from (\ref{InvLDef0}),
that $\rf{\bmi  G}({\bmi x};{\bmi x}^{\prime};k_y)$ 
may be written
\begin{samepage}
\begin{eqnarray} \label{InvLDef}
\rf{\bmi  G}({\bmi x};{\bmi x}^{\prime};k_y)&=&{\CB F}^{-1}\, \left[\rf{\bmi  G}({\bmi 
k};{\bmi x}^{\prime};k_y) \right]({\bmi x}) =\frac{1}{4 \pi^2} 
\int_{R^2 } {\rm d}{\bmi k} \exp \left[i {\bmi k}\cdot ({\bmi x}-{\bmi 
x}^{\prime})\right]\nonumber\\
&&\times {\bmi  D}_{z}({k}_{x},{k}_{y})
\left[{\begin{array}{cc}{\bf 0}_{3}&-{\bmi W}^{+}\\
{\bmi W}^{-}&{\bf 0}_{3}\end{array}}\right]{\bmi  D}_{z}^{T}(-{k}_{x},-{k}_{y})
\,{\bmi  J}_6
\end{eqnarray}
\end{samepage} 

\subsection{The classical up/down, plane wave Green's tensor}
We now evaluate the $k_{z}$ integral in (\ref{InvLDef}) using the calculus of 
residues [e.g. {\em Aki \& Richards\/}~(1980)], which is possible provided
the $z$ component of  $({\bmi x}-{\bmi x}^{\prime})$ in  (\ref{InvLDef}) 
is non-zero, i.e. $({z}-{z}^{\prime})\not=0 $.

We  close the 
$k_z$ integration contour 
in the ${\stackrel{{\mbox{{\small upper }}}}{{}_{\mbox{{\small lower }}}}}$ half of the 
complex ${k}_{z}-$plane when 
$(z-{z}^{\prime})\gtlt0$, in so doing  we pick up the 
\[
\Re(\partial \omega/\partial { k}_{z,c})\gtlt 0
\]
residue contributions (Fig.\ref{Poles}), associated with 
${\stackrel{\mbox{\small downgoing}}{{}_{\mbox{\small upgoing}}}}$ 
group velocity, to get
\begin{eqnarray} \label{Gzres}
\lefteqn{\rf{\bmi  G}({\bmi x};{\bmi x}^{\prime};k_y)=\frac{- i}{2 \pi} 
\int_{R} {\rm d}{ k}_{x} \exp \left[i { k}_{x} ({ x}-{ 
x}^{\prime})\right] {\bmi  D}_{z}({k}_{x},{k}_{y})
{\bmi Q}^{z}(z;z^{\prime}){\bmi  D}_{z}^{T}(-{k}_{x},-{k}_{y})
\,{\bmi  J}_6}\nonumber\\
&&
\end{eqnarray} 
where
${\bmi Q}^{z}(z;z^{\prime})$ in (\ref{Gzres}),
is defined  as
\begin{equation} \label{Qzdef}
{\bmi Q}^{z}(z;z^{\prime})=
\left[{\begin{array}{cc}{\bf 0}_{3}&H(z-{z}^{\prime}){\bmi E}(k_x;z;{z}^{\prime})\\
H(z^{\prime}-{z}){\bmi E}(k_x;z;{z}^{\prime})&{\bf 0}_{3}\end{array}}\right]
\end{equation}
with
\begin{equation} \label{Ek1def}
{\bmi E}(k_x;z;{z}^{\prime})=\left[{\begin{array}{ccc}
{\rm e}^{i{ k}_{z,P}(k_{x})|z-{z}^{\prime}|}&{ 0}&{ 0}\\ 
{ 0}&{\rm e}^{ 
i{ k}_{z,S}(k_{x})|z-{z}^{\prime}|}&{ 0}\\ { 0}&{ 0}&{\rm e}^{ 
i{ k}_{z,H}(k_{x})|z-{z}^{\prime}|}\end{array}}\right] ,
\end{equation}
and $H(\cdot)$ is the {\em Heaviside step function\/}.
The imposition of the {\em radiation boundary conditions\/} employed here, 
follows {\em Lighthill}~\shortcite{Lighthill}. Alternatively,  we may
 introduce {\em attenuation\/} to move the 
poles off the integration contour [e.g. {\em Aki \& Richards\/}~(1980)].


%%%%%%%%%%%%%%%%%%%%%%%%%%%%%%%%%%%%%%%%%%%%%%%%%%%%%%%%%%%%%%%%%%%%%%%%%%%%%%
%%%%%%%%%%%%%%%%%%%%%%%%%%%%%%%%%%%%%%%%%%%%%%%%%%%%%%%%%%%%%%%%%%%%%%%%%%%%%%
%%%%%%%%%%%%%%%%%%%%%%%%%%%%%%%%%%%%%%%%%%%%%%%%%%%%%%%%%%%%%%%%%%%%%%%%%%%%%
\begin{figure}
\HideDisplacementBoxes
%\ShowDisplacementBoxes
\centerline{\BoxedEPSF{Poles.EPSF scaled 800}}
\caption{The movement of poles in the complex $k_z$ plane consistent with
$2 \pi \nu \rightarrow 2 \pi \nu +i \epsilon$}
\protect\label{Poles}
\end{figure}
%%%%%%%%%%%%%%%%%%%%%%%%%%%%%%%%%%%%%%%%%%%%%%%%%%%%%%%%%%%%%%%%%%%%%%%%%%%%%
%%%%%%%%%%%%%%%%%%%%%%%%%%%%%%%%%%%%%%%%%%%%%%%%%%%%%%%%%%%%%%%%%%%%%%%%%%%%%%%%%




\section[Discrete plane wave representations]
{Discrete plane wave representations of the free space 
elastodynamic Green's tensor}\label{secPWG}

\subsection{The $z$ splitting}\label{secStandardz}
\begin{figure}
\fbox{
\begin{minipage} {142.0mm}
{\sc Box} $\bf 3.2$ {\sc The discrete Fourier transform\/}\newline
Given the Fourier integral
\begin{equation} \label{invF}
\phi  (x)={\frac{1}{2 \pi }}\int_{- {\mit \Omega} }^{ {\mit \Omega} }{\rm d}{k}_{x}
\exp \left({i{k}_{x}x}\right)\tilde{ \phi }({k}_{x}).
\end{equation}
we will now show how to approximate it using quadrature.
To begin,  (\ref{invF}) may be approximated by the 
$2M$-panel {\em trapezoidal rule\/} as
\begin{eqnarray} \label{DFT0}
 \phi   (x) &\approx& {\frac{{\mit \Delta} {  k}_{  
x}}{  2 \pi }}\left[{{\frac{  \exp \left({-i {\mit \Omega} 
  x}\right)\tilde{ \phi }(- {\mit \Omega}   )+\exp 
\left({i {\mit \Omega}   x}\right)\tilde{ \phi }( {\mit \Omega}   
)}{2}} }\right] \nonumber\\
&& + \sum\limits_{  j= -M+1}^{  M-1}  
 \exp \left({ij {\mit \Delta} {  k}_{  x}  
x}\right)\tilde{ \phi }(j {\mit \Delta} {  k}_{  x}  
)
\end{eqnarray} 
where the wavenumber sampling interval  ${\mit \Delta} {  k}_{  x}$ is 
defined as
\[
{\mit \Delta} {  k}_{  x}=\frac{2 {\mit \Omega}}{2 M}=\frac{{\mit \Omega}}{M}
\]
For computational purposes, we {\em alias\/} the endpoints in 
(\ref{DFT0}), i.e.
\begin{equation} \label{aliasapprox}
{{\frac{  \exp \left({-i {\mit \Omega} 
  x}\right)\tilde{ \phi }(- {\mit \Omega}   )+\exp 
\left({i {\mit \Omega}   x}\right)\tilde{ \phi }( {\mit \Omega}   
)}{2}} }\equiv 
\exp \left({-i {\mit \Omega} 
  x}\right)\tilde{ \phi }(- {\mit \Omega}   )
\end{equation}
so that (\ref{DFT0}) becomes
\begin{equation} \label{DFT1}
 \phi(x) \approx {\frac{{\mit \Delta} {  k}_{  
x}}{  2 \pi }}\sum\limits_{  j= -M}^{  M-1}  
 \exp \left({ij {\mit \Delta} {  k}_{  x}  
x}\right)\tilde{ \phi }(j {\mit \Delta} {  k}_{  x}  )
\end{equation} 
The approximation (\ref{aliasapprox}) is feasible when
\begin{equation} \label{endpoints}
\tilde{ \phi }(- {\mit \Omega}   )=
\tilde{ \phi }({\mit \Omega}   )= 0
\end{equation}
which is the approximation adopted here.  

We will now proceed to manipulate (\ref{DFT1}) into a {\em discrete Fourier 
transform\/} (DFT). To begin, we rewrite the sum in (\ref{DFT1}) as
\begin{eqnarray} \label{DFT2}
 \phi   (x) &\approx& {\frac{{\mit \Delta} {  k}_{  
x}}{  2 \pi }}\left[{\sum\limits_{  j=0}^{  M-1} 
  \exp \left({ij {\mit \Delta} {  k}_{  x}  
x}\right)\tilde{ \phi }(j {\mit \Delta} {  k}_{  x}  
)+ \sum\limits_{  j=-M}^{  -1}   \exp \left({i j 
{\mit \Delta} {  k}_{  x}  x}\right)\tilde{ \phi }(j 
{\mit \Delta} {  k}_{  x}  )}\right]\nonumber\\
& \approx& {\frac{{\mit \Delta} {  k}_{  x}}{  2 \pi 
}}\left[\sum\limits_{  j=0}^{  M-1}   \exp 
\left({i j {\mit \Delta} {  k}_{  x}  x}\right)\tilde{ 
\phi }(j {\mit \Delta} {  k}_{  x}  ) \right.\nonumber\\
&&\qquad\qquad
\left.+ \sum\limits_{ 
 j=M}^{  2M-1}   \exp \left({i(j-2M) {\mit \Delta} {  
k}_{  x}  x}\right)\tilde{ \phi }((j-2M) {\mit \Delta} { 
 k}_{  x}  )\right]
\end{eqnarray} 
We now discretize $x$ in (\ref{DFT2}) so that
\begin{equation} \label{xgrid}
{x}_{l}=l {\mit {\mit \Delta}}  x,\qquad 0 \le  l \le  2M-1.
\end{equation}

Using  ${x}_{l}$ from (\ref{xgrid}) in place of $x$ in (\ref{DFT2}) we get
\begin{samepage}
\begin{eqnarray} \label{DFT3}
\lefteqn{ \phi   ({x}_{l}) \approx {\frac{{\mit \Delta} {  k}_{ 
 x}}{  2 \pi }}\left[\sum\limits_{  j=0}^{  
M-1}   \exp \left({ijl {\mit \Delta} {  k}_{  x}{\mit \Delta} 
  x}\right)\tilde{ \phi }(j {\mit \Delta} {  k}_{  
x}  )+\exp \left({-i2Ml {\mit \Delta} {  k}_{  x}{\mit \Delta} 
  x}\right)\right.}\nonumber\\
&&
\left.
\qquad
\times \sum\limits_{  j=M}^{  2M-1}   
\exp \left({ijl {\mit \Delta} {  k}_{  x}{\mit \Delta}   
x}\right)\tilde{ \phi }((j-2M) {\mit \Delta} {  k}_{  x}  )\right],
\quad  0 \le l \le 2M-1
\end{eqnarray} 
\end{samepage}

\vspace{3mm}
\end{minipage}
}
\end{figure}

\begin{figure}
\fbox{
\begin{minipage} {142.0mm}
{\sc Box} $\bf 3.2${\em continued\/} \newline


Finally, we impose the fundamental condition in Discrete Fourier analysis 
that
\begin{equation}  \label{Fconstraint}
{\mit \Delta} { k}_{ x}{\mit \Delta}  x={\frac{2 \pi }{ 2M}}
\end{equation}
and employing this constraint in (\ref{DFT3}) we get
\begin{equation}\label{DFT}
 \phi ({x}_{l}) \approx {\frac{{\mit \Delta} { k}_{  x}}{
  2 \pi }}\sum\limits_{ j=0}^{ 2M-1}  \exp \left({{\frac{2 \pi  i j l}{2M}}}\right)
\tilde{ \varphi }(j {\mit \Delta} { k}_{ x} ), \qquad  0 \le l \le 2M-1
\end{equation} 
where we have introduced the {\em DFT-ordered vector\/} $\tilde{\varphi 
}$ defined as
\begin{equation} \label{Forder}
\tilde{\varphi }(j {\mit \Delta} { k}_{ x} )=
\left\{{\begin{array}{cl}\tilde{ \phi }(j {\mit \Delta} { k}_{ x} )&0 \le  j\rm \le  M-1\\ 
\tilde{ \phi 
}([j-2M] {\mit \Delta} { k}_{ x} )&M \le  
j \le  2M-1\end{array}}\right.
\end{equation} 
and  used the identity
\[
\exp \left({-i2 M l  {\mit \Delta} {  k}_{  x}{\mit \Delta}   x}\right)
=\exp \left(-i2 M l  \frac{2 \pi }{ 2M}\right)=1
\]
which follows directly from the constraint (\ref{Fconstraint}). To 
evaluate the DFT (\ref{DFT}) using the {\em fast Fourier transform\/}(FFT) algorithm, 
we impose the further condition 
\[
2M=2^m, \quad m\in\{1,2,3,\cdots  \}
\]
in which case the floating point operation count of (\ref{DFT}) is 
reduced from $O(4 M^2)$ to $O(2 m M)$.

Finally, an efficient method to incorporate {\em Filon quadrature\/} 
into the DFT approximation is discussed in {\sc Box} $\bf 3.3$.
\vspace{3mm}
\end{minipage}
}
\end{figure}

Basic results concerning the discrete Fourier (DFT) approximation to the inverse 
Fourier transform over horizontal wavenumber of a scalar function,
 are summarized in {\sc Box} $\bf 3.2$ and a simple extension of the 
 basic theory to incorporate Filon quadrature, while exploiting the speed 
 of the FFT algorithm is presented in {\sc Box} $\bf 3.3$. 
We will employ these results to develop a DFT approximation to (\ref{Gzres}).  
We assume the $N_k=2 M$ element {\em Fourier lattice\/}  
\begin{equation} \label{pspec}
\left\{ 
 {k}_{j} =j {\mit \Delta} k_{x} \ |\  j=-M,\cdots ,M-1 
\right\} 
\end{equation} 
where ${\mit \Delta} k_{x} =2 \pi  j/L_{x}$ and   $L_{x}$ defines the {\em fundamental spatial period\/} 
of the representation.
We may then approximate (\ref{Gzres}) by the {\em discrete 
plane wave sum\/}
\begin{samepage}
\begin{eqnarray} \label{Gdef5}
\rf{\bmi G}\!{}^{z}({\bmi x};{\bmi x}^{\prime};k_y)&=&\frac{- i\,{\mit \Delta} {k}_{x}}{2 \pi} 
\sum\limits_{{j=-M}}^{j=M-1}
\nonumber\\
&&\times
\exp \left[i {k}_{j} ({ x}-{ 
x}^{\prime})\right] {\bmi  D}_{z}({k}_{j},{k}_{y})
{\bmi Q}^{z}({k}_{j};z;z^{\prime}){\bmi  D}_{z}^{T}(-{k}_{j},-{k}_{y})
\,{\bmi  J}_6
\end{eqnarray} 
\end{samepage}
The discrete representation (\ref{Gdef5}) is equivalent to assuming that the 
transverse dependence of the wavefield is both {\em  bandlimited\/} 
(since $N_k$ is 
finite) and periodically repeated (since ${\mit \Delta} {k}_{x}$ is finite).


\begin{figure}
\fbox{
\begin{minipage} {142.0mm}
{\sc Box} $\bf 3.3$ {\sc An application of systematic Filon 
quadrature to the discrete Fourier transform\/}\newline

An interesting method for {\em attenuating the effects of aliasing\/} in 
the discrete Fourier transform of  {\sc Box} $\bf 3.2$ is to iterpolate 
the spectrum before transforming: a technique commonly employed in the 
Filon method [e.g. {\em Frazer \& Gettrust\/}~(1984), {\em Press \& 
Teukolsky\/}~(1989)]. We will demonstrate this technique with a specific example.

We interpolate the digitized spectrum
\[
\{\tilde{ \phi }(j {\mit \Delta} {  k}_{  x}  ): -M \le j \le M \}
\]
with the {\em chapeau shape function\/} $\tilde{X}(k_x)$ (of finite 
element fame), defined as
\[
\tilde{X}({k}_{x})=\left\{{\begin{array}{cc}1-\left|{{k}_{x}/ {\mit \Delta}   
k_x}\right|&- {\mit \Delta}   k_x \le {  k}_{  x}\le {\mit \Delta} 
  k_x\\ 0&\mbox{otherwise}\end{array}}\right.
\] 
so that the spectrum may be approximated as
\begin{equation} \label{interp}
\tilde{ \Phi }({  k}_{  x}  )=\sum\limits_{j=-M}^{ M} \tilde{ X} 
({k}_{x}-j {\mit \Delta} { k}_{ x})\tilde{ \phi }(j {\mit \Delta} { k}_{ x} ) 
=\tilde{ X}({k}_{x})\, {\ast}
\sum\limits_{j=-M}^{ M} \delta 
({k}_{x}-j {\mit \Delta} { k}_{ x})\tilde{ \phi }(j {\mit \Delta} { k}_{ x} )
\end{equation}
where the binary operator ``$\ast$" denotes a {\em convolution\/} and 
$\delta$ is the Dirac delta function. Applying the inverse Fourier 
transform to  (\ref{interp}) gives
\begin{eqnarray} \label{interp1}
\lefteqn{{\frac{1}{2 \pi }}\int {\rm d}\, {k}_{x}{\rm e}^{i{k}_{x}x}\tilde{ \Phi 
}({k}_{x})={\frac{1}{2 \pi }}\int {\rm d}{k}_{x}{\rm e}^{i{k}_{x}x} 
\sum\limits_{  j=-M}^{  M} \tilde{  X}  
({k}_{x}-j {\mit \Delta} {  k}_{  x}  )\tilde{ \phi 
}(j {\mit \Delta} {  k}_{  x}  )}\nonumber\\
&&\!\!\!\!\!\!
= \left[ {\frac{1}{2 \pi }}\int {\rm d}({k}_{x}-j {\mit \Delta} {  k}_{  x}  )\,
{\rm e}^{i({k}_{x}-j {\mit \Delta} {  k}_{  x}  )x}
 \tilde{  X}  
({k}_{x}-j {\mit \Delta} {  k}_{  x}  )\right]\,\sum\limits_{  j=-M}^{  M} {\rm e}^{i j {\mit \Delta} {  
k}_{  x} x}\,\tilde{ \phi 
}(j {\mit \Delta} {  k}_{  x}  )\nonumber\\
&&\!\!\!\!\!\!
=\left[ {\frac{{\mit \Delta} {  k}_{  x}}{2 \pi }}
\mbox{sinc}^{2}\left( \frac{{\mit \Delta} {  k}_{  x}\, x}{2} \right)
\right]\,\sum\limits_{  j=-M}^{  M} {\rm e}^{i j {\mit \Delta} {  
k}_{  x} x}\,\tilde{ \phi 
}(j {\mit \Delta} {  k}_{  x}  )
\end{eqnarray} 

$\quad$ The extra {\em multiplicative factor\/}
\[
A(x)=\mbox{sinc}^{2}\left(\frac{{\mit \Delta} {  k}_{  x}\, x}{2}\right)
\] 
in (\ref{interp1}) {\em attenuates \/} the basic DFT sum 
\begin{equation} \label{Atendef}
{\frac{{\mit \Delta} {  k}_{  x}}{2 \pi }}\,
\sum\limits_{  j=-M}^{  M} {\rm e}^{i j {\mit \Delta} {  
k}_{  x} x}\,\tilde{ \phi 
}(j {\mit \Delta} {  k}_{  x}  )
\end{equation}
that we met previously in {\sc Box } $3.2$. If the fundamental period 
 of the representation is $L_x$, then the wavenumber sampling interval
 ${\mit \Delta} {  k}_{  x}$ is chosen according to
\[
{\mit \Delta} {  k}_{  x} =\frac{2 \pi }{L_x}
\] 
in which case (\ref{Atendef}) is a {\em monotone decreasing \/}  function of $x$ for $ 0 
\le x \le L_x$. In particular 
\[
A(x)\sim \left\{ \begin{array}{cc}  1 &x \rightarrow 0 \\ 0 & x 
\rightarrow L_x  \end{array}\right.
\]

$\quad$ In this thesis, we always choose $L_x$ so that $0 \le x 
\le L_x/8$, and consequently $A(x)\approx 1$. Even so, if we decide to 
include larger offsets into the basic method, 
the attenuation factor (\ref{Atendef}) could be included as a  
multiplicative correction, at little additional cost. 
\vspace{3mm}
\end{minipage}
}
\end{figure}


%%%%%%%%%%%%%%%%%%%%%%%%%%%%%%%%%%%%%%%%%%%%%%%%%%%%%%%%%%%%%%%%%%%%%%%%%%%%%%%%%%%%%%%%
%%%%%%%%%%%%%%%%%%%%%%%%%%%%%%%%%%%%%%%%%%%%%%%%%%%%%%%%%%%%%%%%%%%%%%%%%%%%%%%%%%%%%%%%
%%%%%%%%%%%%%%%%%%%%%%%%%%%%%%%%%%%%%%%%%%%%%%%%%%%%%%%%%%%%%%%%%%%%%%%%%%%%%%%%%%%%%%%%
The next step will be to rewrite the discrete sum (\ref{Gdef5}) as an 
equivalent matrix operation. To this end, we define the $N_k-$element 
diagonal, 
horizontal-wavevector array:
\begin{equation} \label{gKx}
{\bf  k}_{x} =\mbox{diag}\left[0,
{\mit \Delta} k_{x},\cdots, (M -1){\mit \Delta} k_{x},
-M {\mit \Delta} k_{x},\cdots,-{\mit \Delta} k_{x},\right];
\end{equation}
the $N_k-$element 
diagonal, 
vertical-wavevector array 
\begin{eqnarray} \label{gKz}
\lefteqn{{\bf  k}_{z,c} =\mbox{diag}\left[{ k}_{z,c}(0),{ k}_{z,c}({\mit \Delta} k_{x}),\cdots,
{ k}_{z,c}((M-1){\mit \Delta} k_{x}),
{ k}_{z,c}(-M{\mit \Delta} k_{x}),
\cdots, { k}_{z,c}(-{\mit \Delta} k_{x})\right]}\nonumber\\
&&\hspace{25em} c\in\{P,S,H\};
\end{eqnarray} 
and the $N_k-$element {\em Fourier vector\/}
\begin{equation} \label{Fdef}
{\bmi F}(x;{\bf k}_{x})=
\left[\begin{array}{ccc} 1&, \cdots, &1 \end{array}  
\right]\cdot{\rm e}^{i\, 
{\bf k}_{x}\,x}.
\end{equation} 
It is important to note that ${\bmi F}(x;{\bf k}_{x})$ as defined in 
(\ref{Fdef}) is an $N_k-$ row vector, even though ${\bf k}_{x}$ is an 
$N_k-$diagonal matrix. Also, the elements of (\ref{gKx}),
(\ref{gKz})  and (\ref{Fdef})  have the DFT 
ordering of (\ref{Forder}) for the reasons indicated in {\sc Box} 
$\bf 3.2$.


Using the newly defined arrays (\ref{gKx})--(\ref{Fdef}),
we can  rewrite the {\em discrete plane wave\/} 
approximation (\ref{Gdef5}) to the $6 \times 6$ Green's tensor in the equivalent matrix 
form
\begin{samepage}
\begin{eqnarray} \label{superG}
\rf{\bmi G}\!{}^{z}({\bmi x};{\bmi x}^{\prime};k_y)
&=&- i\,\left({\bmi I}_6\otimes 
{\bmi F}(x;{\bf k}_{x})  \right) {\bmi  D}_{z}({\bf k}_{x},{k}_{y})\nonumber\\
&& \times
{\bmi Q}^{z}({\bf k}_{x};z;z^{\prime}){\bmi  D}_{z}^{T}(-{\bf k}_{x},-{k}_{y})
\,\left({\bmi I}_6\otimes {\bmi F}^{T}(x^{\prime};-{\bf k}_{x})  \right)\,{\bmi  J}_6
\end{eqnarray} 
\end{samepage}
where:\newline
the {\em Kronecka product\/}
$\,\left({\bmi I}_6\otimes {\bmi F}(x^{\prime};{\bf k}_{x})  \right)$ is 
a $  6 \times 6 N_k$ array, 
${\bmi F}^{H}(x;{\bf k}_{x})$ is the {\em Hermitian transpose\/} of the 
Fourier vector ${\bmi F}(x;{\bf k}_{x})$ defined in (\ref{Fdef})
(so that ${\bmi F}^{H}(x;{\bf k}_{x})\equiv {\bmi F}^{T}(x;-{\bf 
k}_{x})$);\newline
 the $(6N_k)\times (6N_k)$ matrix analogue of (\ref{Ddef}) is defined 
as
\begin{eqnarray} \label{Dsuperdef}
\lefteqn{{\bmi D}_{z}({\bf k}_{x},k_{y})=}\\
&&\!\!\!\!\!\!\!\!\!\!\!\!
\left[\begin{array}{cccccc}
\left|+,P, {\bf k}_{x},k_{y}  \right\rangle&\left|+,S, {\bf k}_{x},k_{y}  
\right\rangle&\left|+,H, {\bf k}_{x},k_{y}  \right\rangle&
\left|-,P, {\bf k}_{x},k_{y}  \right\rangle&\left|-,S, {\bf k}_{x},k_{y}  
\right\rangle&\left|-,H, {\bf k}_{x},k_{y}  \right\rangle
\end{array}\right]\nonumber
\end{eqnarray}
in terms of
\begin{itemize}
\item the $6 N_k \times N_k$ $P$-wave partition
\begin{equation} \label{PeigenSuper}
\left|\pm,P, {\bf k}_x,k_{y}  \right\rangle=
\sqrt{\frac{{\mit \Delta} k_{x}}{{4 \pi \rho {\omega 
}^{  2}{\bf  k}_{  
z, P }}}}\,\left[\begin{array}{c} \pm i\,{\bf k}_{z,P}\\i\,{\bf k}_{x}\\
i\,k_{y}\,{\bmi I}\\
  \rho \,\left( 2\,{\beta^2}\,{{\bf k}_{x}^2} + ( 2\,{\beta^2}\,{k_{y}^2} - 
     {\omega^2}){\bmi I} \right) \\ \mp 2\,\rho\,{\beta^2}\,{\bf k}_{x}\,{\bf k}_{z,P} \\
  \mp 2\,\rho\,{\beta^2}\,k_{y}\,{\bf k}_{z,P} \end{array}\right]
\end{equation}

\item  the $6 N_k \times N_k$ $SV$-wave partition
\begin{equation} \label{SVeigenSuper}
\left|\pm,S, {\bf k}_x,k_{y}  \right\rangle=\frac{\omega }{\beta \hat{K}_{S}}\,
\sqrt{\frac{{\mit \Delta} k_{x}}{{4 \pi \rho {\omega 
}^{  2}{\bf  k}_{  
z, S }}}}\,\left[\begin{array}{c} 
i\,{\bf k}_{x}\\
\mp i\,{\bf k}_{z,S}\\ {\bf 0} \\
\mp 2\,\rho\,{\beta^2}\,{\bf k}_{x}\,{\bf k}_{z,S} \\
  \rho \,\left(({\omega^2} - {\beta^2}\,{k_{y}^2}) {\bmi I}-2\,{\beta^2}\,{{\bf k}_{x}^2}   \right) \\
  -\rho\,{\beta^2}\,k_{y}\,{\bf k}_{x} \end{array}\right]
\end{equation}

\begin{samepage}
\item the $6 N_k \times N_k$ $SH$-wave partition
\begin{equation} \label{SHeigenSuper}
\left|\pm,H, {\bf k}_x, k_{y}  \right\rangle=\frac{1 }{\hat{K}_{H}}\,
\sqrt{\frac{{\mit \Delta} k_{x}}{{4 \pi \rho {\omega 
}^{  2}{\bf  k}_{  
z, H }}}}\,\left[\begin{array}{c} 
\mp k_{y}\,{\bf k}_{z,H} \\-k_{y}\,{\bf k}_{x} \\
  \hat{K}_{H}^2 {\bmi I}\\
-2\,i\,\rho \,{\beta^2}\,k_{y}\,{{\bf k}_{z,H}^2}\\
  \mp 2\,i\,\rho\,{\beta^2}\,k_{y}\,{\bf k}_{x}\,{\bf k}_{z,H} \\
  \pm i\,\rho\,\left(  {\omega^2} -2\,{\beta^2}\,{k_{y}^2} 
  \right)\,{\bf k}_{z,H}  \end{array}\right]
\end{equation}
\end{samepage}

\end{itemize}
The $(6N_k)\times (6N_k)$ {\em plane wave propagator\/} 
${\bmi Q}^{z}({\bf k}_{x};z;z^{\prime})$ in (\ref{superG}),
is defined  as
\begin{equation} \label{Qdef}
{\bmi Q}^{z}({\bf k}_x;z;z^{\prime})=
\left[{\begin{array}{cc}{\bf 0}_{3N_k}&H(z-{z}^{\prime}){\bmi E}({\bf k}_x;z;{z}^{\prime})\\
H(z^{\prime}-{z}){\bmi E}({\bf k}_x;z;{z}^{\prime})&{\bf 0}_{3N_k}\end{array}}\right]
\end{equation}
where the $(3N_k)\times (3N_k)$ {\em vertical phase shift array\/} is 
defined as
\[
{\bmi E}({\bf k}_x;z;{z}^{\prime})=\left[{\begin{array}{ccc}
{\rm e}^{i{\bf k}_{z,P}|z-{z}^{\prime}|}&{\bf 0}_{N_k}&{\bf 0}_{N_k}\\ 
{\bf 0}_{N_k}&{\rm e}^{ 
i{\bf k}_{z,S}|z-{z}^{\prime}|}&{\bf 0}_{N_k}\\ {\bf 0}_{N_k}&{\bf 0}_{N_k}&{\rm e}^{ 
i{\bf k}_{z,H}|z-{z}^{\prime}|}\end{array}}\right] ,
\]
in terms of the diagonal wavevector array ${\bf k}_{z,c}, c\in\{P,S,H  
\}$ of (\ref{gKz}).

It will be convenient in later applications,  to split (\ref{superG}) into the product
\[
\rf{\bmi G}\!{}^{z}({\bmi x};{\bmi x}^{\prime};k_y)=
\rf{\bmi P}\!{}^{z}({\bmi x};{\bmi x}^{\prime};k_y)\cdot{\bmi J}_6
\]
where ${\bmi J}_6$ was given by (\ref{Jdef}) and
\begin{samepage}
\begin{eqnarray} \label{superP}
\rf{\bmi P}\!{}^{z}({\bmi x};{\bmi x}^{\prime};k_y)
&=&- i\,\left({\bmi I}_6\otimes 
{\bmi F}(x;{\bf k}_{x})  \right) {\bmi  D}_{z}({\bf k}_{x},{k}_{y})\nonumber\\
&& \times
{\bmi Q}^{z}({\bf k}_{x};z;z^{\prime}){\bmi  D}_{z}^{T}(-{\bf k}_{x},-{k}_{y})
\,\left({\bmi I}_6\otimes {\bmi F}^{T}(x^{\prime};-{\bf k}_{x})  \right).
\end{eqnarray} 
\end{samepage}
might be called the {\em propagator\/}, for the want of a better name. The 
propagator is an entity in itself for most numerical applications 
considered here.

The representation (\ref{superG}) (or equivalently, (\ref{superP})) is readily extensible 
to the plane stratified problem and we investigate this extension in 
chapter {\ref{stratG}}. 

\subsection{The $x$ splitting }\label{secxsplit}


Recalling the definition (\ref{Ddef}) for the columns of the eigen-matrix 
${\bmi  D}_{z}$, (\ref{InvLDef}) may be rewritten as 
\begin{samepage} 
\begin{eqnarray}\label{grep2}
\rf{\bmi  G}({\bmi  x};{{\bmi  x}}^{\prime};k_y)& =&\sum\limits_{c\in \{P,S,H\}}
{\frac{1}{4{ \pi }^{2}}}\int_{{R}^{2}} {\rm d}{\bmi k} \exp \left[{i {\bmi k}  \cdot 
  ({\bmi x}-{\bmi x}^{\prime})}\right]\nonumber\\
&&\!\!\!\!\!\!\!\!\!\!\!\!\!\!\!\!\!\!
\times\left[{{\frac{
  \left|-,c, k_{x},k_{y}  \right\rangle\otimes \left|+,c, -k_{x},-k_{y}  \right\rangle
  }{{k}_{z}+{k}_{z, c }({k}_{x})}}-{\frac{
\left|+,c, k_{x},k_{y}  \right\rangle\otimes \left|-,c, -k_{x},-k_{y}  \right\rangle
  }{{k}_{z}-{k}_{z, c }({k}_{x})}}}\right]{\bmi J}_{6}
\nonumber\\
&=&
\rf{\bmi  G}_{P}({\bmi  x};{{\bmi  x}}^{\prime};k_y)+
\rf{\bmi  G}_{S}({\bmi  x};{{\bmi  x}}^{\prime};k_y)+
\rf{\bmi  G}_{H}({\bmi  x};{{\bmi  x}}^{\prime};k_y)
\end{eqnarray} 
\end{samepage}
where  $\otimes$ is the {\em tensor product\/}.
We will ultimately evaluate the integrals in (\ref{grep2}) using the 
residues at the simple poles, where  
\begin{equation} \label{dispersion}
{k}^2_{z}+{k}^2_{x}+{k}^2_{y}=\frac{\omega^2}{c^2},\quad \mbox{for}\quad c\in\{P,S,H\}.
\end{equation}
Hence,  (\ref{grep2}) may be replaced by an expression that is 
{\em equivalent at these poles\/}. Introducing the partial differential 
operators
\begin{eqnarray} \label{Labdef}
\begin{array}{lll}
{\bmi L}_{P}=&\qquad{\bmi L}_{S}=&\qquad{\bmi L}_{H}=\\
\left[\begin{array}{c}
\partial_{z} \\
\partial_{x} \\ 
\partial_{y} \\ 
2 \mu \partial_{zz} - {\lambda \omega^{2}\over \alpha^2}\\
    2 \mu \partial_{xz} \\ 
   2 \mu \partial_{yz}
\end{array}\right]&\qquad
\frac{\omega }{\beta \hat{K}_{S}}\,
\left[\begin{array}{c}
\partial_{x} \\-\partial_{z} \\ 0 \\
2 \mu \partial_{xz}\\  
   \mu (-\partial_{zz} + \partial_{xx}) \\ 
   \mu \partial_{xy} 
\end{array}\right]&\qquad
\frac{1 }{\hat{K}_{H}}\,
\left[\begin{array}{c}
   \partial_{yz} \\
\partial_{xy} \\ 
{\omega^{2}\over \beta^2} + \partial_{yy} \\ 
 2 \mu \partial_{yzz}\\
    2 \mu \partial_{xyz} \\ 
   \mu ({\omega^{2}\over \beta^2}\partial_{z} + 2 \partial_{yyz}) 
\end{array}\right]
\end{array}
\end{eqnarray}
$\rf{\bmi G}_{c}({\bmi  x};{{\bmi  x}}^{\prime};k_y)$  in  (\ref{grep2}) may be replaced 
with the {\em equivalent\/} expression
\begin{eqnarray}\label{grepc}
\rf{\bmi G}_{c}({\bmi  x};{{\bmi  x}}^{\prime};k_y) &=&
{\frac{1}{8{ \pi }^{2} \rho \omega^2}}\int_{{R}^{2}} {\rm d}{\bmi k}\,  
\frac{
{\bmi L}_{c} \otimes {\bmi  L}_{c}^{\prime}
\left(
\exp \left[{i {\bmi k}  \cdot  ({\bmi x}-{\bmi x}^{\prime})}\right]\right)  
}{ {k}_{z, c }({k}_{x})\,[{k}_{z}+{k}_{z, c }({k}_{x})]} {\bmi J}_{6}\nonumber\\
&&-
{\frac{1}{8  { \pi }^{2}\rho \omega^2}}\int_{{R}^{2}} {\rm d}{\bmi k}\,  
\frac{{\bmi L}_{c} \otimes {\bmi  L}_{c}^{\prime}\left(
\exp \left[{i {\bmi k}  \cdot  ({\bmi x}-{\bmi x}^{\prime})}\right]\right)  
}{{k}_{z, c }({k}_{x})\, [{k}_{z}-{k}_{z, c }({k}_{x})]} {\bmi J}_{6},\  
c\in\{ P, S , H\};\nonumber\\
&=& \frac{1}{4{ \pi }^{2}\rho \omega^2}
\int_{{R}^{2}} {\rm d}{\bmi k}\, {\bmi L}_{c} \otimes {\bmi  L}_{c}^{\prime}
\frac{
\exp \left[{i {\bmi k}  \cdot  ({\bmi x}-{\bmi x}^{\prime})}\right]  
}{ \hat{K}^{2}_{ c } - {k}^{2}_{z} - {k}^{2}_{x}} {\bmi J}_{6}
\end{eqnarray} 
where we have used the algebraic relation (\ref{dispersion}) and
 the differential operator ${\bmi  L}_{c}^{\prime}$ operates on the
${\bmi  x}^{\prime}$ coordinate. 


Assuming that $x \not = x^{\prime}$ in (\ref{grepc}) we may  evaluate the $k_x$ integral in  
that equation using the calculus of residues to get
\begin{samepage}
\begin{eqnarray} \label{grep3x}
\lefteqn{\rf{\bmi G}_{c}({\bmi  x};{{\bmi  x}}^{\prime};k_y)=
\frac{-i}{2{ \pi }\rho \omega^2}
\int_{{R}} {\rm d}{k}_z \,{\bmi L}_{c } \otimes {\bmi  L}_{c}^{\prime}
\frac{\exp \left[{i {k}_z  ({z}-{z}^{\prime})}\right]
\exp \left[{i {k}_{x,c}({k}_z,{k}_y)  |{x}-{x}^{\prime}|}\right]  
}{ 2\,{k}_{x,c}({k}_z,{k}_y)}{\bmi J}_{6},}\nonumber\\
&&
\end{eqnarray}
\end{samepage}
with 
\begin{equation}	\label{kxcDef}
	{k}_{x,c}({k}_{z},{k}_{y}) = \left\{
	\begin{array}{c}
	\sqrt {\hat{ K}_{c}^{2}({k}_{y})\ -\ {k}_{z}^{2}}\ \ \ \ \ \ |{k}_{x}|\ <\ 
	|\hat{K}_{c}({k}_{y})|\\
	i\sqrt {{k}_{z}^{ 2}\ -\ \hat{ K}_{c}^{2}({k}_{y})} \ \ \ \ |{k}_{x}|\ \ge   \ 
	|\hat{K}_{c}({k}_{y})|
\end{array}\right.
\end{equation}
and $\hat{ K}_{c}({k}_{y})$ was given by (\ref{Kdef}).
Using (\ref{grep3x}) in (\ref{grep2}) 
we get
\begin{eqnarray} \label{GxGRep}
\rf{\bmi G}_{c}({\bmi  x};{{\bmi  x}}^{\prime};k_y)&=&
\frac{-i}{2{ \pi }\rho \omega^2}\,\sum\limits_{c\in \{P,S,H\}}\nonumber\\
&\times& 
\int_{{R}} {\rm d}{k}_z \,{\bmi L}_{c } \otimes {\bmi  L}_{c}^{\prime}
\frac{\exp \left[{i {k}_z  ({z}-{z}^{\prime})}\right]
\exp \left[{i {k}_{x,c}({k}_z,{k}_y)  |{x}-{x}^{\prime}|}\right]  
}{ 2\,{k}_{x,c}({k}_z,{k}_y)}{\bmi J}_{6}
\end{eqnarray}

Using the definition of the 
partial differential operators ${\bmi L}_{c }$ from (\ref{Labdef}) in (\ref{grep2}) and 
performing the relevant spatial differentiations  we get
\begin{eqnarray} \label{Gxres}
\lefteqn{\rf{\bmi G}({\bmi x};{\bmi x}^{\prime};k_y)=\frac{- i}{2 \pi} 
\int_{R} {\rm d}{ k}_{z} \exp \left[i { k}_{z} ({ z}-{ 
z}^{\prime})\right] {\bmi  D}_{x}({k}_{z},{k}_{y})
{\bmi Q}^{x}(k_z;x;x^{\prime}){\bmi  D}_{x}^{T}(-{k}_{z},-{k}_{y})
\,{\bmi  J}_6}\nonumber\\
&&
\end{eqnarray} 
provided $x \not = x^{\prime}$, where:\newline
\begin{eqnarray} \label{Dxdef}
\lefteqn{{\bmi D}_{x}(k_{z},k_{y})=}\\
&&\!\!\!\!\!\!\!\!\!\!\!\!
\left[\begin{array}{cccccc}
\overline{\left|+,P, k_{z},k_{y}  \right\rangle}&\overline{\left|+,S, k_{z},k_{y}  
\right\rangle}&\overline{\left|+,H, k_{z},k_{y}  \right\rangle}&
\overline{\left|-,P, k_{z},k_{y}  \right\rangle}&\overline{\left|-,S, k_{z},k_{y}  
\right\rangle}&\overline{\left|-,H, k_{z},k_{y}  \right\rangle}
\end{array}\right]\nonumber
\end{eqnarray}
where:
\begin{itemize}
\item the $\overline{\left|\pm,P, k_{z},k_{y}  \right\rangle}$ states 
are the   ${\stackrel{\mbox{\small rightgoing}}{{}_{\mbox{\small 
leftgoing}}}}$  $P-$waves 
\begin{equation} \label{Pxeigen}
\overline{\left|\pm,P, k_{x,P},k_{y}  \right\rangle}=\frac{1}{\sqrt{2 \rho {\omega }^{  2}{  k}_{  
x, P }}}\,\left[\begin{array}{c}  
i\,{k}_{z}\\ \pm i\,k_{x,P}(k_x,k_y)\\i\,k_{y}\\
  \rho \,\left( 2\,{\beta^2}\,{k_{x,P}^2(k_x,k_y)} + 2\,{\beta^2}\,{k_{y}^2} - 
     {\omega^2} \right) \\ \mp 2\,\rho\,{\beta^2}\,k_{x,P}(k_x,k_y)\,{k}_{z} \\
  - 2\,\rho\,{\beta^2}\,k_{y}\,{k}_{z} \end{array}\right]
\end{equation}

\begin{samepage}
\item the $\overline{\left|\pm,S, k_{z},k_{y}  \right\rangle}$ states 
are the   ${\stackrel{\mbox{\small rightgoing}}{{}_{\mbox{\small 
leftgoing}}}}$  quasi-$SV-$waves 
\begin{equation} \label{SVxeigen}
\overline{\left|\pm,S, k_{z},k_{y}  \right\rangle}=\frac{\omega}{\beta\,\hat{K}_{S }\,
\sqrt{2 \rho {\omega }^{  2}  {  k}_{  
x, S }}}\,
\left[\begin{array}{c} 
\pm i\,k_{x,S}(k_x,k_y)\\
- i\,{k}_{z}\\0\\
\mp 2\,\rho\,{\beta^2}\,k_{x,S}(k_x,k_y)\,{k}_{z} \\
  \rho \,\left({\omega^2} -2\,{\beta^2}\,{k_{x,S}^2(k_x,k_y)} - {\beta^2}\,{k_{y}^2}  \right) \\
  \mp \rho\,{\beta^2}\,k_{x,S}(k_x,k_y)\,k_{y} \end{array}\right]
\end{equation}
\end{samepage}

\item the $\overline{\left|\pm,H, k_{z},k_{y}  \right\rangle}$ states 
are the   ${\stackrel{\mbox{\small rightgoing}}{{}_{\mbox{\small 
leftgoing}}}}$  quasi-$SH-$waves 
\begin{equation} \label{SHxeigen}
\overline{\left|\pm,H, k_{z},k_{y}  \right\rangle}=
\frac{1}{\hat{K}_{H }\, \sqrt{2 \rho {\omega }^{  2}  {  k}_{  
x, H }}} \,\left[\begin{array}{c} 
- k_{y}\,{k}_{z} \\\mp k_{x,H}(k_x,k_y)\,k_{y} \\
  \hat{K}_{H}^2\\
2\,i\,k_{y}\,\rho \,\left({\beta^2}\,{k_{x,H}^2(k_x,k_y)} + 
     {\beta^2}\,{k_{y}^2} - {\omega^2} \right) \\
  \mp  2\,i\,\rho\,{\beta^2}\,k_{x,H}(k_x,k_y)\,k_{y}\,{k}_{z}\\
   i\,{k}_{z}\,\rho \,\left(  {\omega^2} -2\,{\beta^2}\,{k_{y}^2} \right) \end{array}\right]
\end{equation}
\end{itemize}
the {\em horizontal plane wave propagator\/} ${\bmi Q}^{x}(k_z;x;x^{\prime})$ in (\ref{Gxres}),
is defined  as
\begin{equation} \label{Qxdef}
{\bmi Q}^{x}(k_z;x;x^{\prime})=
\left[{\begin{array}{cc}{\bf 0}_{3}&H(x-{x}^{\prime}){\bmi E}(k_z;x;{x}^{\prime})\\
H(x^{\prime}-{x}){\bmi E}(k_z;x;{x}^{\prime})&{\bf 0}_{3}\end{array}}\right]
\end{equation}
with the {\em horizontal phase shift\/} array
\[
{\bmi E}(k_z;x;{x}^{\prime})=\left[{\begin{array}{ccc}
{\rm e}^{i{ k}_{x,P}(k_{z})|x-{x}^{\prime}|}&{ 0}&{ 0}\\ 
{ 0}&{\rm e}^{ 
i{ k}_{x,S}(k_{z})|x-{x}^{\prime}|}&{ 0}\\ { 0}&{ 0}&{\rm e}^{ 
i{ k}_{x,H}(k_{z})|x-{x}^{\prime}|}\end{array}}\right] ,
\]
and $H(\cdot)$ is the {\em Heaviside step function\/}.


%%%%%%%%%%%%%%%%%%%%%%%%%%%%%%%%%%%%%%%%%%%%%%%%%%%%%%%%%%%%%%%%%%%%%%%%%%%%%%%%%%%%%%%%%%
%%%%%%%%%%%%%%%%%%%%%%%%%%%%%%%%%%%%%%%%%%%%%%%%%%%%%%%%%%%%%%%%%%%%%%%%%%%%%%%%%%%%%%%%%%
%%%%%%%%%%%%%%%%%%%%%%%%%%%%%%%%%%%%%%%%%%%%%%%%%%%%%%%%%%%%%%%%%%%%%%%%%%%%%%%%%%%%%%%%%%
For computational purposes, we must employ a {\em discrete plane wave\/} 
or {\em Fourier lattice\/} version of the wavenumber integral 
(\ref{Gxres}). The development closely parallels that of \S\ref{secStandardz}.
One of the many problems associated with discrete Fourier representations 
of the Green's tensor operators commonly employed seismology is 
wraparound. In the standard reflectivity method, we employ a horizontal 
wavenumber sampling, and as a consequence we get aliased versions of the 
wavefield {\em wrapping\/} into the computational domain from the sides 
(Mallick \& Frazier, 1988). This is not such a serious problem provided 
we extend the transverse support of the computational domain, or 
equivalently, decrease the horizontal wavenumber sampling interval 
interval   ${\mit \Delta} k_{x}$.

We will now derive a {\em discrete plane wave\/} version of 
(\ref{Gxres}), and we may expect wraparound in $z$. This is a far more 
serious problem than wraparound in $x$ because we will eventually extend 
the representation to $z-$stratified media, so that the wraparound phases 
will interact with the stratification in a complicated manner, that may 
not be removed.
However, we 
will employ (\ref{Gxres}) only in the special case where $z=z^{\prime}$ 
so that $\exp \left[i \overline{k}_{j} ({ z}-{ 
z}^{\prime})\right]\equiv 1$ and consequently 
\begin{eqnarray} \label{Gxres0}
\lefteqn{\rf{\bmi G}({ x};{ x}^{\prime};k_y)=\frac{- i}{2 \pi} 
\int_{R} {\rm d}{ k}_{z}  {\bmi  D}_{x}({k}_{z},{k}_{y})
{\bmi Q}^{x}(k_z;x;x^{\prime}){\bmi  D}_{x}^{T}(-{k}_{z},-{k}_{y})
\,{\bmi  J}_6, \ z=z^{\prime}, \ x\not=x^{\prime} }\nonumber\\
&&
\end{eqnarray} 
so there will be no {\em 
aliasing\/} or {\em wraparound\/} in the $z-$coordinate. This observation 
is crucial to the success in the intended application of $\rf{\bmi G}({\bmi 
x};{\bmi x}^{\prime};k_y)$ from (\ref{Gxres0}) to problems with a 
$z-$stratification in material properties. 


Given the $\overline{N}_k=2 \overline{M}+1$ element {\em Fourier lattice\/}  
\begin{equation} \label{pxspec}
\left\{ 
 \overline{k}_{j} =j {\mit \Delta} k_{z} \ |\  j=-\overline{M},\cdots 
 ,\overline{M} 
\right\} 
\end{equation} 
where ${\mit \Delta} k_{z} =2 \pi  j/L_{z}$ and   $L_{z}$ 
controls the  {\em wavenumber sampling rate,\/} we may approximate 
(\ref{Gxres0}) by the {\em discrete 
plane wave sum\/}
\begin{equation} \label{Gxdef5}
\rf{\bmi G}\!{}^{x}({ x};{ x}^{\prime};k_y)=\frac{- i\,{\mit \Delta} {k}_{z}}{2 \pi} 
\sum\limits_{{j=-\overline{M}}}^{j=\overline{M}}
{\bmi  D}_{x}(\overline{k}_{j},{k}_{y})
{\bmi Q}^{x}(\overline{k}_{j};x;x^{\prime}){\bmi  D}_{x}^{T}(-\overline{k}_{j},-{k}_{y})
\,{\bmi  J}_6
\end{equation} 
The discrete representation (\ref{Gxdef5}) is equivalent to assuming that the 
$z-$dependence of the wavefield is {\em  bandlimited\/} (since $\overline{M}$ is 
finite), which is certainly not the case in the limit $x \leftarrow 
x^{\prime}$.



%%%%%%%%%%%%%%%%%%%%%%%%%%%%%%%%%%%%%%%%%%%%%%%%%%%%%%%%%%%%%%%%%%%%%%%%%%%%%%%%%%%%%%%%
%%%%%%%%%%%%%%%%%%%%%%%%%%%%%%%%%%%%%%%%%%%%%%%%%%%%%%%%%%%%%%%%%%%%%%%%%%%%%%%%%%%%%%%%
%%%%%%%%%%%%%%%%%%%%%%%%%%%%%%%%%%%%%%%%%%%%%%%%%%%%%%%%%%%%%%%%%%%%%%%%%%%%%%%%%%%%%%%%
The next step will be to rewrite the discrete sum (\ref{Gxdef5}) as an 
equivalent matrix operation. To this end, we define the diagonal wave 
vector arrays:
\begin{equation} \label{gbarKz}
{\bf  k}_{z} =\mbox{diag}\left[{-\overline{M} {\mit \Delta} k_{z},\cdots,-{\mit 
\Delta} k_{z},0,
{\mit \Delta} k_{z},\cdots, \overline{M} {\mit \Delta} k_{z}}\right];
\end{equation} 
\begin{equation} \label{gbarKx}
{\bf  k}_{x,c} =\mbox{diag}\left[{{ k}_{x,c}(-\overline{M} {\mit \Delta} k_{z}),\cdots,
{ k}_{x,c}(-{\mit \Delta} k_{z}),
{ k}_{x,c}(0),
{ k}_{z,c}({\mit \Delta} k_{z}),\cdots, { k}_{x,c}(\overline{M}{\mit \Delta} k_{z})}\right]
\end{equation} 
and the $\overline{N}_k-$element {\em projection \/} vector
\begin{equation} \label{Fbardef}
\overline{\bmi F}(0)=\left[\begin{array}{ccccc}
1& \cdots & 
1& \cdots & 
1\end{array}\right].
\end{equation} 


We can  rewrite the {\em discrete plane wave\/} 
approximation (\ref{Gxdef5}) to the $6 \times 6$ Green's tensor in the equivalent matrix 
form
\begin{samepage}
\begin{eqnarray} \label{superGx}
\rf{\bmi G}\!{}^{x}({ x};{ x}^{\prime};k_y)
&=&- i\, 
\left({\bmi I}_6\otimes \overline{\bmi F}(0)  \right)
 {\bmi  D}_{x}({\bf k}_{z},{k}_{y})\nonumber\\
&&\times
{\bmi Q}^{x}({\bf k}_{z};x;x^{\prime}){\bmi  D}_{x}^{T}(-{\bf k}_{z},-{k}_{y})
\,\left({\bmi I}_6\otimes \overline{\bmi F}^{T}(0)  \right)\,{\bmi  J}_6
\end{eqnarray} 
\end{samepage}
where:\newline
the $(6\bar{N}_k)\times (6\bar{N}_k)$ matrix analogue of (\ref{Dxdef}) is defined 
as
\begin{eqnarray} \label{Dxsuperdef}
\lefteqn{{\bmi D}_{x}({\bf k}_{z},k_{y})=}\\
&&\!\!\!\!\!\!\!\!\!\!\!\!
\left[\begin{array}{cccccc}
\overline{\left|+,P, {\bf k}_{z},k_{y}  \right\rangle}&\overline{\left|+,S, {\bf k}_{z},k_{y}  
\right\rangle}&\overline{\left|+,H, {\bf k}_{z},k_{y}  \right\rangle}&
\overline{\left|-,P, {\bf k}_{z},k_{y}  \right\rangle}&\overline{\left|-,S, {\bf k}_{z},k_{y}  
\right\rangle}&\overline{\left|-,H, {\bf k}_{z},k_{y}  \right\rangle}
\end{array}\right]\nonumber
\end{eqnarray}
in terms of
\begin{itemize}
\item the $6 \overline{N}_k \times \overline{N}_k$ $P$-wave partition
\begin{equation} \label{PxeigenSuper}
\overline{\left|\pm,P, {\bf k}_z,k_{y}  \right\rangle}=
\sqrt{\frac{{\mit \Delta} k_{z}}{{4 \pi \rho {\omega 
}^{  2}{\bf  k}_{  
x, P }}}}
\left[\begin{array}{c} i\,{\bf k}_{z}\\ \pm i\,{\bf k}_{x,P}\\
i\,k_{y}\,{\bmi I}\\
  \rho \,\left( 2\,{\beta^2}\,{{\bf k}_{x,P}^2} + ( 2\,{\beta^2}\,{k_{y}^2} - 
     {\omega^2}){\bmi I} \right) \\ \mp 2\,\rho\,{\beta^2}\,{\bf k}_{x,P}\,{\bf k}_{z} \\
  - 2\,\rho\,{\beta^2}\,k_{y}\,{\bf k}_{z} \end{array}\right]
\end{equation}

\item  the $6 \overline{N}_k \times \overline{N}_k$ $SV$-wave
\begin{equation} \label{SVxeigenSuper}
\overline{\left|\pm,S, {\bf k}_z,k_{y}  \right\rangle}=
\frac{\omega }{\beta \hat{K}_{S}}\,\sqrt{\frac{{\mit \Delta} k_{z}}{{4 \pi \rho {\omega 
}^{  2}{\bf  k}_{  
x, S }}}}\,\left[\begin{array}{c} 
\pm i\,{\bf k}_{x,S}\\
- i\,{\bf k}_{z}\\ {\bf 0} \\
\mp 2\,\rho\,{\beta^2}\,{\bf k}_{x,S}\,{\bf k}_{z} \\
  \rho \,\left(({\omega^2} - {\beta^2}\,{k_{y}^2}) {\bmi 
  I}-2\,{\beta^2}\,{{\bf k}_{x,S}^2}   \right) \\
 \mp \rho\,{\beta^2}\,k_{y}\,{\bf k}_{x,S} \end{array}\right]
\end{equation}

\item the $6 \overline{N}_k \times \overline{N}_k$ $SH$-wave partition
\begin{equation} \label{SHxeigenSuper}
\overline{\left|\pm,H, {\bf k}_z, k_{y}  \right\rangle}=\frac{1 }{\hat{K}_{H}}\,
\sqrt{\frac{{\mit \Delta} k_{z}}{{4 \pi \rho {\omega 
}^{  2}{\bf  k}_{  
x, H }}}}\,\left[\begin{array}{c} 
- k_{y}\,{\bf k}_{z} \\ \mp k_{y}\,{\bf k}_{x,H} \\
  \hat{K}_{H}^2 {\bmi I}\\
-2\,i\,\rho \,{\beta^2}\,k_{y}\,{{\bf k}_{z}^2}\\
  \mp 2\,i\,\rho\,{\beta^2}\,k_{y}\,{\bf k}_{x,H}\,{\bf k}_{z} \\
   i\,\rho\,\left(  {\omega^2} -2\,{\beta^2}\,{k_{y}^2} 
  \right)\,{\bf k}_{z}  \end{array}\right]
\end{equation}

\end{itemize}
where  ${\bf  k}_{x}$ and ${\bf  k}_{z,c}$ were given in
(\ref{gbarKx}) and (\ref{gbarKz}) respectively.

The $(6\bar{N}_k)\times (6\bar{N}_k)$ {\em plane wave propagator\/} 
${\bmi Q}^{x}({\bf k}_{z};x;x^{\prime})$ in (\ref{superGx}),
is defined  as
\begin{equation} \label{gQxdef}
{\bmi Q}^{x}({\bf k}_z;x;x^{\prime})=
\left[{\begin{array}{cc}{\bf 0}_{3\bar{N}_k}&H(x-{x}^{\prime}){\bmi E}({\bf 
k}_z;x;{x}^{\prime})\\
H(x^{\prime}-{x}){\bmi E}({\bf k}_z;x;{x}^{\prime})&{\bf 0}_{3\bar{N}_k}\end{array}}\right]
\end{equation}
where the $(3\bar{N}_k)\times (3\bar{N}_k)$ {\em vertical phase shift array\/} is 
defined as
\[
{\bmi E}({\bf k}_z;x;{x}^{\prime})=\left[{\begin{array}{ccc}
{\rm e}^{i{\bf k}_{z,P}|x-{x}^{\prime}|}&{\bf 0}_{\overline{N}_k}&{\bf 0}_{\overline{N}_k}\\ 
{\bf 0}_{\overline{N}_k}&{\rm e}^{ 
i{\bf k}_{z,S}|x-{x}^{\prime}|}&{\bf 0}_{\overline{N}_k}\\ {\bf 0}_{\overline{N}_k}&{\bf 0}_{\overline{N}_k}&
{\rm e}^{ 
i{\bf k}_{z,H}|x-{x}^{\prime}|}\end{array}}\right] ,
\]
in terms of the diagonal wavevector array ${\bf k}_{x,c}, c\in\{P,S,H  
\}$ of (\ref{gbarKx}).

For computational efficiency, we split (\ref{superGx}) into
\[
\rf{\bmi G}\!{}^{x}({ x};{ x}^{\prime};k_y)=\rf{\bmi P}\!{}^{x}({ x};{ x}^{\prime};k_y){\bmi  J}_6
\]
where
\begin{samepage}
\begin{eqnarray} \label{superPx}
\rf{\bmi P}\!{}^{x}({ x};{ x}^{\prime};k_y)
&=&- i\, 
\left({\bmi I}_6\otimes \overline{\bmi F}(0)  \right)
 {\bmi  D}_{x}({\bf k}_{z},{k}_{y})\nonumber\\
&&\times
{\bmi Q}^{x}({\bf k}_{z};x;x^{\prime}){\bmi  D}_{x}^{T}(-{\bf k}_{z},-{k}_{y})
\,\left({\bmi I}_6\otimes \overline{\bmi F}^{T}(0)  \right)
\end{eqnarray} 
\end{samepage}




\section{Lattice propagator representations in free space} \label{latticeP}
This section is central to the thesis. In this section the medium is 
homogeneous. The method to be expounded here is extended to a stratified 
halfspace in chapter \ref{stratG}. 

To begin, we define a  uniform lattice (Fig.\ref{lattice})
with $L$ rows and ${N}_{x}$ lattice 
points per row.
The lattice points are numbered in the 
same order that words are read on the written page, 
i.e. from left to right, from top to bottom.
 For many applications it is helpful to label the lattice 
points as if they were the elements of a $2-$D array. For this purpose we 
introduce the matrix row index $1\le \rho \le L$ and the matrix column 
index $1 \le \chi \le N_x$, so that the global index $i$ is related to 
the matrix indices via
\[
i=(\rho -1) N_x + \chi, \quad 1 \le i \le L N_x, 
\quad 1 \le \rho \le L, \quad 1 \le \chi \le N_x
\] 
Please note that use of $\rho$ and $\chi$ to indicate the row and column 
indices respectively, are confined to this section, and are not to 
be confused with the other applications of these symbols throughout the thesis.


We want to {\em 
communicate\/} the influence of the {\em secondary force field\/} $6-$vector
 $\bar{\bmi f}_{ j}$, from a given lattice point 
to all the {\em other\/} lattice points.
Using the naive matrix/vector operation
\begin{eqnarray} \label{GFdef0}
\rf{\CB G} \cdot \bar{\CB F} &\equiv& \sum\limits_{ \stackrel{j= 1}{j\not= i}}^{ 
L{N}_{x}} \rf{\bmi G}({\bmi R}_{i};{\bmi R}_{j};{k}_{y}) \cdot \bar{\bmi f}_{ j}
\qquad 1\le  i \le L{N}_{x}
\end{eqnarray} 
where   $\rf{\bmi G}$ is the $6 \times 6$ Green's 
tensor (\ref{(2.1.6)}), (\ref{GFdef0}) may be implemented as
\[
N_1=(6 \times {N}_{x} \times 
L)^2 - (6 \times {N}_{x} \times L)
\]
separate interactions.


%%%%%%%%%%%%%%%%%%%%%%%%%%%%%%%%%%%%%%%%%%%%%%%%%%%%%%%%%%%%%%%%%%%%%%%%%%%%%
\begin{figure}
\HideDisplacementBoxes
%\ShowDisplacementBoxes
\centerline{\BoxedEPSF{lattice.EPSF scaled 800}}
\caption{The lattice geometry}
\protect\label{lattice}
\end{figure}
%%%%%%%%%%%%%%%%%%%%%%%%%%%%%%%%%%%%%%%%%%%%%%%%%%%%%%%%%%%%%%%%%%%%%%%%%%%%%

We will employ the plane wave expansions (\ref{superG}), (\ref{superGx}),
[or equivalently,  (\ref{superP}), (\ref{superPx})] , the  Fast Fourier 
Transform and 
matrix factorization techniques so that (\ref{GFdef0}) is computed using
\begin{eqnarray} \label{GFdef}
\rf{\CB G} \cdot \bar{\CB F} &\equiv& \sum\limits_{ \stackrel{j= 1}{j\not= i}}^{ 
L{N}_{x}} \left[
(1-\delta_{\rho \rho^{\prime}})
\rf{\bmi P}\!{}^{z}({\bmi R}_{i};{\bmi R}_{j};{k}_{y}) \cdot {\bmi J}_6 \bar{\bmi f}_{ j}
+
\delta_{\rho \rho^{\prime}}\rf{\bmi P}\!{}^{x}({\bmi R}_{i};{\bmi R}_{j};{k}_{y}) 
 \cdot {\bmi J}_6 \bar{\bmi f}_{ j}\right]\nonumber\\
&\equiv&- i\,\sum\limits_{ \stackrel{j= 1}{j\not= i}}^{ L{N}_{x}}
(1-\delta_{\rho \rho^{\prime}})\left({\bmi I}_6\otimes 
{\bmi F}(X_i;{\bf k}_{x})  \right) {\bmi  D}_{z}({\bf k}_{x},{k}_{y})\nonumber\\
&& \times 
{\bmi Q}^{z}({\bf k}_{x};Z_i;Z_j){\bmi  D}_{z}^{T}(-{\bf k}_{x},-{k}_{y})
\,\left({\bmi I}_6\otimes {\bmi F}^{T}(x^{\prime};-{\bf k}_{x})  \right)
\,\cdot\,{\bmi f}_{ j}\nonumber\\
&&- i\, \sum\limits_{ \stackrel{j= 1}{j\not= i}}^{ L{N}_{x}}
\delta_{\rho \rho^{\prime}}
\left({\bmi I}_6\otimes \overline{\bmi F}(0)  \right)
 {\bmi  D}_{x}({\bf k}_{z},{k}_{y})\nonumber\\
&&\times
 {\bmi Q}^{x}({\bf k}_{z};X_i;X_j){\bmi  D}_{x}^{T}(-{\bf k}_{z},-{k}_{y})
\,\left({\bmi I}_6\otimes \overline{\bmi F}^{T}(0)  \right)\,\cdot\,{\bmi f}_{ j}\nonumber\\
&&\qquad i\in \mbox{row}_{\rho}, \qquad j\in \mbox{row}_{\rho^{\prime}}, \qquad 
\mbox{for}\ \mbox{each} \quad  1\le  i \le L{N}_{x}
\end{eqnarray} 
where
\[
{\bmi f}_{ j}\equiv {\bmi J}_6 \bar{\bmi f}_{ j}
\]
to reduce the  floating point operation count of (\ref{GFdef0}) to
\[
N_2\approx2\,(L-1)\,\left[ 39 N_k + 6 N_k \log_2 N_k \right]+ 
78\,L\,(N_x-1)  \overline{N}_k
\]
where $N_k$ is the number of plane waves employed to propagate the 
secondary source field between rows, while $\bar{N}_k$ is the number of plane
 waves employed to propagate the 
secondary source field between lattice points within a given row.
If we employ critical (i.e. Nyquist) sampling rates so that $N_k=N_x$, then $N_2 \ll N_1$. 
However, for the sake of accuracy, we need to employ a substantial amount 
of oversampling so that $N_k>N_x$, which increases the operation count 
$N_2$ of our plane wave expansion method. In the applications presented 
in chapter \ref{variational}, $N_k=16 N_x$ or even $N_k=32 N_x$, so that the plane 
wave method is not substantially faster than the naive operation count 
$N_1$; at least for the scale of the problems tested, for larger problems 
we observe that $N_2$ grows much more slowly than $N_1$. The principal merit of the method to 
be presented in this section lies in its 
extensibility to a plane stratified  medium.

Recall that we have introduced a so called {\em propagator\/}, denoted as 
$\rf{\bmi P}\!{}^{z,x}({\bmi 
R}_{i};{\bmi R}_{j};k_y)$, and defined by reference to the Green's 
tensors of (\ref{superG}) and (\ref{superGx}). We will denote the generic 
propagator as $\rf{\bmi P}\!({\bmi 
R}_{i};{\bmi R}_{j};k_y)$, and will re-introduce the superscripts in 
$\rf{\bmi P}\!\!{}^{z}({\bmi 
R}_{i};{\bmi R}_{j};k_y)$ and $\rf{\bmi P}\!\!{}^{x}({\bmi 
R}_{i};{\bmi R}_{j};k_y)$, when appropriate.
Let us assemble these $6 \times 6$ propagators $\rf{\bmi P}\!\!({\bmi 
R}_{i};{\bmi R}_{j};k_y)$  for lattice points ${\bmi R}_{i}$ and ${\bmi 
R}_{j}$ into a $6 {N}_{x} L \times 6  {N}_{x} L$ {\em super-matrix\/} 
${\CB P}$ 
where the $(i, j)$ partition ${\CB P}_{i,j}$ of ${\CB P}$ is defined as
\begin{equation} \label{Gijdef}
{\CB P}_{i,j}=(1-{\delta }_{ij})\rf{\bmi P}({\bmi R}_{i};{\bmi 
R}_{j};k_y)
\end{equation}
so that  (\ref{GFdef}) may be expressed as the equivalent matrix vector operation
\begin{equation} \label{CBGdef}
{\CB P}\cdot{\CB F}=\left[{\begin{array}{cccc}
{\bf 0}_6&{\CB P}_{1,2}&\cdots &{\CB P}_{1,{N}_{x} L}\\
{\CB P}_{2,1}&{\bf 0}_6&&\vdots\\
\vdots&&\ddots&{\CB P}_{{N}_{x} L-1,{N}_{x} L}\\
{\CB P}_{{N}_{x} L,1}&\cdots &{\CB P}_{{N}_{x} L,{N}_{x} L-1}&{\bf 0}_6\end{array}}\right]
\cdot\left[\begin{array}{c} {\bmi f}_1\\
{\bmi f}_2\\\vdots\\{\bmi f}_{L N_x}    \end{array}\right]
\end{equation}

We assemble each of the $6-$vectors ${\bmi f}_{{N}_{x}({\rho 
}^{\prime}-1)+\chi}, \ 1 \le \chi \le N_x$ in $\mbox{row}_{\rho^{\prime}}$ of the 
lattice into the $6 N_x$-vector
\begin{equation} \label{Frho}
{\CB F}_{{\rho }^{\prime}}=\left[{\begin{array}{c}
{\bmi f}_{{N}_{x}({\rho }^{\prime}-1)+1}\\
\vdots\\
{\bmi f}_{{N}_{x}{\rho }^{\prime}}\end{array}}\right]
\end{equation}
From (\ref{GFdef}), we observe  that the propagator partitions $\rf{\bmi P}({\bmi 
R}_{i};{\bmi R}_{j})$ have different plane wave 
expansions depending on whether the lattice points ${\bmi 
R}_{i}$ and ${\bmi R}_{j}$ are within the same or different rows of the 
lattice. Consequently, we label the partitions of (\ref{CBGdef}) depending 
on the source and receiver row, i.e.
\begin{equation} \label{CBdef1}
{\CB P}=\left[{\begin{array}{cccc}{\CB P}_{1 \leftarrow   1}&{\CB P}_{1 
\leftarrow   2}&\cdots &{\CB P}_{1 \leftarrow   L}\\ {\CB P}_{2 
\leftarrow   1}&{\CB P}_{2 \leftarrow   2}&&\vdots\\ 
\vdots&&\ddots&{\CB P}_{L-1 \leftarrow   L}\\ {\CB P}_{L \leftarrow   
1}&\cdots &{\CB P}_{L \leftarrow   L-1}&{\CB P}_{L \leftarrow   
L}\end{array}}\right] 
\end{equation}
where the $6 {N}_{x} \times 6 {N}_{x}$ partition ${\CB P}_{\rho 
\leftarrow   \rho^{\prime}}$ defines 
the {\em influence\/} of the source field ${\CB F}_{\rho^{\prime}}$ 
(\ref{Frho}), 
from the ${N}_{x}$ lattice points in 
$\mbox{row}_{\rho^{\prime}}$ on the ${N}_{x}$ lattice points of 
$\mbox{row}_{\rho}$,  and is defined as
\begin{eqnarray} \label{CBdef2}
\lefteqn{{\CB P}_{\rho \leftarrow   \rho^{\prime}}=\left[{\begin{array}{cccc}
{\CB P}_{(\rho-1) {N}_{x}+1,(\rho^{\prime}-1) {N}_{x}+1}&{\CB P}_{(\rho-1) {N}_{x}+1,(\rho^{\prime}-1) {N}_{x}+2}&\cdots &{\CB P}_{(\rho-1) {N}_{x}+1,\rho^{\prime} {N}_{x}}\\
{\CB P}_{(\rho-1) {N}_{x}+2,(\rho^{\prime}-1) {N}_{x}+1}&{\CB P}_{(\rho-1) {N}_{x}+ 2,(\rho^{\prime}-1) {N}_{x}+2}&&\vdots\\
\vdots&&\ddots&{\CB P}_{\rho {N}_{x}-1 , \rho^{\prime} {N}_{x}}\\
{\CB P}_{\rho {N}_{x},(\rho^{\prime}-1) {N}_{x}+1}&\cdots &{\CB P}_{\rho {N}_{x},\rho^{\prime} {N}_{x}-1}&{\CB P}_{\rho {N}_{x},\rho^{\prime} {N}_{x}}
\end{array}}\right]}\nonumber\\
&&
\end{eqnarray}
$1 \le \rho,\rho^{\prime} \le L$, with $6\times 6$ component arrays ${\CB P}_{i,j}$ 
given by (\ref{Gijdef}). We now split (\ref{CBdef1}) into its off-diagonal 
and diagonal blocks to get
\begin{equation} \label{CBGsplit}
{\CB P}={\CB P}{}^{z} + {\CB P}{}^{x}
\end{equation}
where
\begin{equation} \label{CBdef3}
{\CB P}{}^{z}=\left[{\begin{array}{cccc}
{\bf 0}_{6N_x}&{\CB P}_{1 
\leftarrow   2}&\cdots &{\CB P}_{1 \leftarrow   L}\\ {\CB P}_{2 
\leftarrow   1}&{\bf 0}_{6N_x}&&\vdots\\ 
\vdots&&\ddots&{\CB P}_{L-1 \leftarrow   L}\\ {\CB P}_{L \leftarrow   
1}&\cdots &{\CB P}_{L \leftarrow   L-1}&{\bf 0}_{6N_x}\end{array}}\right] 
\end{equation}
defines the interaction between {\em distinct\/} rows of the lattice
and
\begin{equation} \label{CBdef4}
{\CB P}^{x}=\mbox{diag}\left[\begin{array}{cccc}
{\CB P}_{1 \leftarrow   1},&{\CB P}_{2 
\leftarrow   2},&\cdots, &{\CB P}_{L \leftarrow   L}
\end{array}\right] 
\end{equation}
defines the interaction between points within the  {\em same\/} row of the 
lattice.


Our next task will be to generalize the definition of (\ref{superG}), (\ref{superGx})
 to find a computational construction for (\ref{CBdef3}) and 
(\ref{CBdef4}), respectively.
%%%%%%%%%%%%%%%%%%%%%%%%%%%%%%%%%%%%%%%%%%%%%%%%%%%%%%%%%%%%%%%%%%%%%%%%%%%%%%%%%%%%%%%%
%%%%%%%%%%%%%%%%%%%%%%%%%%%%%%%%%%%%%%%%%%%%%%%%%%%%%%%%%%%%%%%%%%%%%%%%%%%%%%%%%%%%%%%%
%%%%%%%%%%%%%%%%%%%%%%%%%%%%%%%%%%%%%%%%%%%%%%%%%%%%%%%%%%%%%%%%%%%%%%%%%%%%%%%%%%%%%%%%

\subsection{ Implementing $\mbox{\CBM  P}\,{}^{z}\cdot\mbox{\CBM  F}$}

 
 
 
 
Using $\rf{\bmi P}{}^{z}({\bmi x};{\bmi x}^{\prime};k_y)$ from 
(\ref{superP}) for the partitions of
${\CB P}_{\rho \leftarrow   \rho^{\prime}}$ in 
(\ref{CBdef2}), we get
\begin{eqnarray} \label{CBGrhorho1}
{\CB P}_{ \rho \leftarrow {\rho }^{\prime}}
&=&-i\left[{\begin{array}{c}
\left({\bmi I}_6\otimes 
{\bmi F}(x_1;{\bf k}_{x})  \right)\\ \vdots\\ \left({\bmi I}_6\otimes 
{\bmi F}({x}_{{N}_{x}};{\bf k}_{x})\right)\end{array}}\right]
{\bmi  D}_{z}({\bf k}_{x},{k}_{y}) {\bmi Q}_{ \rho 
\leftarrow {\rho }^{\prime}}^{z}({\bf k}_{x})
{\bmi  D}^{T}_{z}(-{\bf k}_{x},-{k}_{y})\nonumber\\
&&\times \left[{\begin{array}{ccc}
\left({\bmi I}_6\otimes{\bmi F}^{T}(x_1;-{\bf k}_{x})  \right)&\cdots &
\left({\bmi I}_6\otimes 
{\bmi F}^{T}({x}_{{N}_{x}};-{\bf k}_{x})\right)\end{array}}\right], 
\quad \rho \not= {\rho }^{\prime}
\end{eqnarray} 
where
\begin{equation} \label{Qdefg}
{\bmi Q}^{z}_{\rho\leftarrow {\rho}^{\prime}}({\bf k}_{x})=
\left[{\begin{array}{cc}{\bf 0}_{3N}&H(z_{\rho}-z_{{\rho}^{\prime}}){\bf E}({\bf k}_{x})^{|\rho-{\rho}^{\prime}|}\\
H(z_{{\rho}^{\prime}}-z_{\rho}){\bf E}({\bf k}_{x})^{|\rho-{\rho}^{\prime}|}&{\bf 0}_{3N}\end{array}}\right]
\end{equation}
with
\[
{\bf E}({\bf k}_{x})=\left[{\begin{array}{ccc}
{\rm e}^{i{\bf k}_{z,P}\,{\mit \Delta}z}&{\bf 0}_{N}&{\bf 0}_{N}\\ 
{\bf 0}_{N}&{\rm e}^{ 
i{\bf k}_{z,S}\,{\mit \Delta}z}&{\bf 0}_{N}\\ {\bf 0}_{N}&{\bf 0}_{N}&{\rm e}^{ 
i{\bf k}_{z,H}\,{\mit \Delta}z}\end{array}}\right] ,
\]
${\mit \Delta}z$ the vertical separation between adjacent rows of the 
lattice in Fig.\ref{lattice}, $H(\cdot)$ is the {\em Heaviside step function\/} 
and the diagonal wavevector array ${\bf k}_{z,c}, c\in\{P,S,H  \}$ was 
given by (\ref{gKz}) and is an implicit function of ${\bf k}_{x}$.


%%%%%%%%%%%%%%%%%%%%%%%%%%%%%%%%%%%%%%%%%%%%%%%%%%%%%%%%%%%%%%%%%%%%%%%%%%%%%%%%%%%
%%%%%%%%%%%%%%%%%%%%%%%%%%%%%%%%%%%%%%%%%%%%%%%%%%%%%%%%%%%%%%%%%%%%%%%%%%%%%%%%%%%
We will now show how to evaluate the matrix/vector product 
\[
{\CB P}_{ \rho \leftarrow {\rho }^{\prime}}\cdot{\CB F}_{{\rho }^{\prime}}
\]
using the fast Fourier transform (FFT) algorithm. To utilize the FFT 
algorithm, we must  transform the $6 N_x$-vector
 ${\CB F}_{{\rho }^{\prime}}$ from (\ref{Frho}) into
the $6 N_k$-vector $\tilde{\CB F}_{{\rho }^{\prime}}\equiv 
{\bmi \Phi}[{\CB F}_{{\rho }^{\prime}}]$ where
\begin{eqnarray} \label{upsample}
{\bmi \Phi}\left[{\CB F}_{{\rho }^{\prime}}\right]\left({N_k (n-1)+m\,{N}_{\mit \Delta }}\right)
&=&\left\{
\begin{array}{cl}
 {\CB F}_{{\rho }^{\prime}}\left({6(m-1)+n}\right)& 1 \le  m \le N_x, \quad 1 \le  n \le 6\\
 0& \mbox{otherwise}
 \end{array}\right.\nonumber\\
 &&
\end{eqnarray} 
with 
\begin{equation} \label{Nkdef}
N_k = {N}_{\mit \Delta } {N}_{L } N_x.
\end{equation}
The mapping (\ref{upsample}) and the corresponding inverse mapping are 
summarized in the pseudo-code of Alg.3.1.


We observe that (\ref{upsample}) defines a dual purpose mapping:
\begin{itemize}
\item firstly, we re-organize the components of ${\CB F}_{{\rho 
}^{\prime}}$ defined in (\ref{Frho}). We recall that all $6$ components 
of the force vector ${\bmi f}_{{N}_{x}({\rho }^{\prime}-1)+\chi}$ for 
each  $ 1 \le 
\chi \le N_x$ are stored contiguously in the $6 N_x$-vector ${\CB 
F}_{{\rho }^{\prime}}$. To employ the DFT, we need to store each fixed 
component, at all gridpoints in a contiguous fashion. This may be 
explained in another way,  ${\CB F}_{{\rho 
}^{\prime}}$ is constructed with an outer loop over each of the $N_x$ the lattice points and 
an inner loop over each of the $6$-components, while the order of the 
inner and outer loop has been swapped in the construction of ${\bmi \Phi}[{\CB F}_{{\rho 
}^{\prime}}]$, so that we have an outer loop over the $6$-components and 
an inner loop over the lattice points.
\item secondly, the mapping ${\bmi \Phi}[\cdot]$ defines an upsampling, 
due to the ${N}_{\mit \Delta }$ factor in (\ref{upsample}). The reason 
for the upsampling will be explained presently.
\end{itemize}

Also the integer relation (\ref{Nkdef}) requires further explanation.  While the width of the 
lattice is  $x_{N_x}\equiv   N_x\,{\mit \Delta }x$, the corresponding {\em Nyquist 
sampling rate\/} for the $k_x$ wavenumber based on this width, i.e.
\[
{\mit \Delta}k_x=\frac{2 \pi }{ N_x\,{\mit \Delta }x},
\]
is usually too coarse, resulting in a significant amount aliasing.  
To get a finer value of  ${\mit \Delta}k_x$, we extend the width 
of the computational grid by a factor ${N}_{L }$ (Fig.\ref{xlattice}) 
so that the actual value 
of ${\mit \Delta}k_x$ is given by
\[
{\mit \Delta}k_x=\frac{2 \pi }{N_L\, N_x\,{\mit \Delta }x}
\]
Also, the wavenumber band width $-{\mit \Omega} \le k_x \le {\mit 
\Omega}$ is usually very wide, (${\mit \Omega}\approx 20 \omega/\beta $ 
in the applications of chapter \ref{variational}) so that the  {\em Nyquist 
spatial sampling rate\/} value of ${\mit \Delta }x$, defined as
\[
{\mit \Delta }x=\frac{ \pi }{{\mit \Omega}} 
\]
is much finer than the actual lattice spacing ${\mit \Delta }x$ in 
Fig.\ref{lattice} and Fig.\ref{xlattice}. Consequently, we have introduced the 
positive integer ${N}_{\mit \Delta }$ in (\ref{Nkdef}) so that the 
lattice spacing is a multiple of the corresponding Nyquist spatial sampling rate 
value, i.e.
\[
{\mit \Delta }x=\frac{{N}_{\mit \Delta } \pi }{{\mit \Omega}} 
\]
The positive integers ${N}_{\mit \Delta }$ and  ${N}_{L }$ can be chosen 
to improve the accuracy of the plane wave representation (\ref{CBGrhorho1}). 
However, this  improvement comes at the cost of employing ${N}_{\mit 
\Delta }\times {N}_{L}$ more wavenumbers in (\ref{CBGrhorho1}) than are 
required by {\em Shannon's sampling theorem\/}.

%%%%%%%%%%%%%%%%%%%%%%%%%%%%%%%%%%%%%%%%%%%%%%%%%%%%%%%%%%%%%%%%%%%%%%%%%%%%%%%%%%%
%%%%%%%%%%%%%%%%%%%%%%%%%%%%%%%%%%%%%%%%%%%%%%%%%%%%%%%%%%%%%%%%%%%%%%%%%%%%%
\begin{figure}
\HideDisplacementBoxes
%\ShowDisplacementBoxes
\centerline{\BoxedEPSF{xlattice.EPSF scaled 800}}
\caption{The $x$-lattice superimposed on the computational grid.
 Circles are the lattice points and 
bullets are the corresponding grid points at the  Nyquist sampling rate. 
$L_x$ is the fundamental 
period of the Fourier representation, while $x_{N_x}$ is the width of the 
lattice}
\protect\label{xlattice}
\end{figure}
%%%%%%%%%%%%%%%%%%%%%%%%%%%%%%%%%%%%%%%%%%%%%%%%%%%%%%%%%%%%%%%%%%%%%%%%%%%%%
%%%%%%%%%%%%%%%%%%%%%%%%%%%%%%%%%%%%%%%%%%%%%%%%%%%%%%%%%%%%%%%%%%%%%%%%%%%%%%%%%%%

The upsampled version of ${\CB P}_{ \rho \leftarrow {\rho }^{\prime}}
\cdot
{\CB F}_{{\rho }^{\prime}}$ from
 (\ref{CBGrhorho1}), corresponding the 
oversampled grid (a single row of that grid is depicted by the ``bullets" 
in Fig.\ref{xlattice}) and operating on the $6 N_k-$element upsampled 
force vector ${\bmi \Phi}\left[{\CB F}_{{\rho }^{\prime}}\right]$ of 
(\ref{upsample}), and may be written as
\[
\tilde{\CB P}_{ \rho \leftarrow {\rho }^{\prime}}\cdot{\bmi \Phi}\left[
{\CB F}_{{\rho }^{\prime}}\right]
\] 
where the $6 N_k \times 6 N_k $ {\em upsampled Green's tensor\/} $\tilde{\CB P}_{ \rho 
\leftarrow {\rho }^{\prime}}$ is an augmented version of $6 N_x \times 6 
N_x $ ${\CB P}_{ \rho 
\leftarrow {\rho }^{\prime}}$ from (\ref{CBGrhorho1}) and is defined as
\begin{samepage}
\begin{eqnarray} \label{tildeCBGrhorho1}
\lefteqn{\tilde{\CB P}_{ \rho \leftarrow {\rho }^{\prime}}
=}\\
&&-i \left({\bmi I}_6\otimes 
{\bf DFT}({\bf k}_x) \right)
{\bmi  D}_{z}({\bf k}_{x},{k}_{y}) {\bmi Q}_{ \rho 
\leftarrow {\rho }^{\prime}}^{z}({\bf k}_{x})
{\bmi  D}^{T}_{z}(-{\bf k}_{x},-{k}_{y})\left({\bmi I}_6\otimes
{\bf DFT}^{T}(-{\bf k}_x)  \right)\nonumber
\end{eqnarray} 
\end{samepage}
In (\ref{tildeCBGrhorho1}), the  ${\bf DFT}({\bf k}_{x})$ represents the $N_k \times N_k$ 
discrete Fourier transform array, with elements
\begin{equation} \label{DFTdef}
\left[{\bf DFT}({\bf k}_x)\right]_{l,j}=\exp \left({{\frac{2 \pi  i (l-1)(j-1)}{N_k}}}\right), \qquad 1 
\le l,j \le N_k   
\end{equation}
while  the elements of {\em Hermitian transpose\/} ${\bf DFT}^{T}(-{\bf 
k}_x)$   
 are defined as
\begin{eqnarray} \label{DFTHdef}
\lefteqn{\left[{\bf DFT}^{T}(-{\bf k}_x)\right]_{l,j}=
\left[{\bf DFT}^{H}({\bf k}_x)\right]_{l,j}=
\exp \left({{\frac{-2 \pi  i (l-1)(j-1)}{N_k}}}\right), \quad 1 
\le l,j \le N_k}\nonumber\\
&&   
\end{eqnarray}
We recall that the Kronecka product $\left({\bmi I}_6\otimes 
{\bf DFT}({\bf k}_x) \right)$ in (\ref{tildeCBGrhorho1})  may be written more explicitly as
the $6 N_k \times 6 N_k$ block diagonal array
\begin{samepage}
\begin{eqnarray*}
\lefteqn{\left({\bmi I}_6\otimes {\bf DFT}({\bf k}_x) \right)\equiv}\\
&&
\mbox{diag}\left[{\begin{array}{cccccc}
{\bf DFT}({\bf k}_x)&{\bf DFT}({\bf k}_x)&{\bf DFT}({\bf k}_x)&
{\bf DFT}({\bf k}_x)&{\bf DFT}({\bf k}_x)&{\bf DFT}({\bf k}_x)\end{array}}\right]
\end{eqnarray*}
\end{samepage}
%%%%%%%%%%%%%%%%%%%%%%%%%%%%%%%%%%%%%%%%%%%%%%%%%%%%%%%%%%%%%%%%%%%%%%%%%%%%%%%%%%%%%
%%%%%%%%%%%%%%%%%%%%%%%%%%%%%%%%%%%%%%%%%%%%%%%%%%%%%%%%%%%%%%%%%%%%%%%%%%%%%%%%%%%%%%%%%%%
\begin{samepage}
\begin{figure}
\fbox{
\begin{minipage} {140.0mm}
{\bf  Algorithm}~{\bf 3.1} {\sc The upsampling ${\bmi \Phi}[\cdot]$ and downsampling 
${\bmi \Phi}^{-1}[\cdot]$ operators}

\begin{tt}

Complex $\tilde{\CB F}_{\rho}(N_k,6), {\CB F}_{\rho}(6,N_x)$\newline
C0=(0,0)


\begin{itemize}
\item Initialize  $\tilde{\CB F}_{\rho}$


\begin{itemize}
\item For j=1,6
\begin{itemize}
\item For k=1,$N_k$
\begin{itemize}
\item $\tilde{\CB F}_{\rho}$(k,j)=C0
\end{itemize}
\item End !k-loop
\end{itemize}
\item End !j-loop
\end{itemize}

\vspace{1 mm}

\item The upsampling  mapping $\tilde{\CB F}_{\rho}={\bmi \Phi}[{\CB F}_{\rho}]$

\begin{itemize}
\item For j=1,6
\begin{itemize}
\item For k=1,$N_x$
\begin{itemize}
\item $\tilde{\CB F}_{\rho}$((k-1)$N_{\mit 
\Delta}+$1,j)=${\CB F}_{\rho}$(j,k)
\end{itemize}
\item End !k-loop
\end{itemize}
\item End !j-loop
\end{itemize}



\end{itemize}

\vspace{2 mm}

\begin{itemize}
\item The downsampling  mapping ${\CB F}_{\rho}={\bmi \Phi}^{-1}[\tilde{\CB F}_{\rho}]$
\begin{itemize}
\item For j=1,6
\begin{itemize}
\item For k=1,$N_x$
\begin{itemize}
\item ${\CB F}_{\rho}$(j,k)$=\tilde{\CB F}_{\rho}$((k-1)$N_{\mit 
\Delta}+$1,j)
\end{itemize}
\item End !k-loop
\end{itemize}
\item End !j-loop
\end{itemize}
\end{itemize}

\end{tt} 

\vspace{2 mm}
\end{minipage}
}
\end{figure}
\vfill
\end{samepage}
%%%%%%%%%%%%%%%%%%%%%%%%%%%%%%%%%%%%%%%%%%%%%%%%%%%%%%%%%%%%%%%%%%%%%%%%%%%%%%%%%%%%%%%%%%%
%%%%%%%%%%%%%%%%%%%%%%%%%%%%%%%%%%%%%%%%%%%%%%%%%%%%%%%%%%%%%%%%%%%%%%%%%%%%%%%%%%%%%

To summarize, to implement (\ref{CBGrhorho1}) using the DFT algorithm, we:
\begin{itemize}
\item {\em upsample\/} the $6 N_x$-element lattice force vector ${\CB F}_{{\rho }^{\prime}}$ 
from (\ref{Frho}) using the 
upsampling operator ${\bmi \Phi}[\cdot]$ of (\ref{upsample}) to get the 
$6 N_k$-element computational vector $\tilde{\CB F}_{{\rho }^{\prime}}$

\item we perform the  mapping 
\[
\tilde{\CB P}_{ \rho \leftarrow {\rho }^{\prime}}\cdot{\bmi \Phi}[{\CB F}_{{\rho 
}^{\prime}}],
\] 
where the $6 N_k \times 6 N_k$ Green's tensor $\tilde{\CB P}_{ \rho 
\leftarrow {\rho }^{\prime}}$ mapping the lattice force on 
$\mbox{row}_{\rho^{\prime}}$ to the 
incident displacement/traction  field on upsampled $\mbox{row}_{\rho}$ 
({\em cf.\/} Fig.\ref{xlattice}),  was defined 
in (\ref{tildeCBGrhorho1}).

\item we apply the inverse operation ${\bmi \Phi}^{-1}[\cdot]$ to the 
result, i.e.
\[
{\bmi \Phi}^{-1}\left[\tilde{\CB P}_{ \rho \leftarrow {\rho }^{\prime}}
\cdot{\bmi \Phi}[{\CB F}_{{\rho 
}^{\prime}}]\right]
\]
that down-samples or {\em decimates\/} the field, to the {\em circles\/} 
in Fig.\ref{xlattice}, while re-ordering the components, so that they are 
stored with an outer loop over offset, and an inner loop over the 
$6-$components of displacement an traction.
\end{itemize}

If we discount the cost of the transformations ${\bmi \Phi}[\cdot]$, 
${\bmi \Phi}^{-1}[\cdot]$, then the floating point operation count of 
$\tilde{\CB P}_{ \rho \leftarrow {\rho }^{\prime}}
\cdot{\bmi \Phi}[{\CB F}_{{\rho 
}^{\prime}}]$ is 
approximately
\[
N_3\approx 75\,{N}_{k}+12\ {N}_{k}^{2}
\]
However, provided $N_k= 2^n$ for some positive integer $n$,
the DFT operations in (\ref{tildeCBGrhorho1}) may be performed with the 
FFT algorithm, in which case $\tilde{\CB P}_{ \rho \leftarrow {\rho }^{\prime}}
\cdot{\bmi \Phi}[{\CB F}_{{\rho 
}^{\prime}}]$ may be implemented in
\[
N_4 \approx 75\,{N}_{k}+12\ {N}_{k} \log_2 N_k
\]
floating point operations.


Next, we consider the mapping
\begin{equation} \label{CBdef03}
{\CB P}{}^{z}\cdot {\CB F}=\left[{\begin{array}{cccc}
{\bf 0}_{6N_x}&{\CB P}_{1 
\leftarrow   2}&\cdots &{\CB P}_{1 \leftarrow   L}\\ {\CB P}_{2 
\leftarrow   1}&{\bf 0}_{6N_x}&&\vdots\\ 
\vdots&&\ddots&{\CB P}_{L-1 \leftarrow   L}\\ {\CB P}_{L \leftarrow   
1}&\cdots &{\CB P}_{L \leftarrow   L-1}&{\bf 0}_{6N_x}\end{array}}\right] 
\left[\begin{array}{c}
{\CB F}_{1}\\ {\CB F}_{2}\\
\vdots\\
{\CB F}_{L}  \end{array}  \right]
\end{equation}
where ${\CB P}{}^{z}$ was introduced in (\ref{CBGsplit}),  the $6 N_x \times 6 N_x$ partition ${\CB P}_{\rho 
\leftarrow   \rho^{\prime}}$ was defined previously by (\ref{CBGrhorho1})
and the $6 N_x$ secondary force vector ${\CB F}_{\rho^{\prime}}$
was introduced in  (\ref{CBGdef}). Once again, we will utilize the DFT (or 
more particularly, the FFT) algorithm to implement (\ref{CBdef03}) in a 
$3$ step process:
\begin{itemize}
\item we introduce the upsampled lattice force vector $\tilde{\CB F}$, 
defined as
\[
\tilde{\CB F}=
\left[\begin{array}{c}
{\bmi \Phi}[{\CB F}_{1}]\\ {\bmi \Phi}[{\CB F}_{2}]\\
\vdots\\
{\bmi \Phi}[{\CB F}_{L}]  \end{array}  \right]
=
\left[\begin{array}{c}
\tilde{\CB F}_{1}\\ \tilde{\CB F}_{2}\\
\vdots\\
\tilde{\CB F}_{L}  \end{array}  \right]
\]
\item we replace  (\ref{CBdef03})  by the upsampled equation
\begin{equation} \label{tildeCBdef3}
\tilde{\CB P}{}^{z}\cdot \tilde{\CB F}=\left[{\begin{array}{cccc}
{\bf 0}_{6N_x}&\tilde{\CB P}_{1 
\leftarrow   2}&\cdots &\tilde{\CB P}_{1 \leftarrow   L}\\ \tilde{\CB P}_{2 
\leftarrow   1}&{\bf 0}_{6N_x}&&\vdots\\ 
\vdots&&\ddots&\tilde{\CB P}_{L-1 \leftarrow   L}\\ \tilde{\CB P}_{L \leftarrow   
1}&\cdots &\tilde{\CB P}_{L \leftarrow   L-1}&{\bf 0}_{6N_x}\end{array}}\right] 
\left[\begin{array}{c}
\tilde{\CB F}_{1}\\ \tilde{\CB F}_{2}\\
\vdots\\
\tilde{\CB F}_{L}  \end{array}  \right]
\end{equation}
where the $6 N_k \times 6 N_k$ partition $\tilde{\CB P}_{\rho 
\leftarrow   \rho^{\prime}}$ was defined previously by 
(\ref{tildeCBGrhorho1}).

\item We then downsample the relevant blocks of $\tilde{\CB P}{}^{z}\cdot 
\tilde{\CB F}$ using the down sampling operator ${\bmi \Phi}^{-1}[\cdot]$
\end{itemize}

However, $\tilde{\CB P}{}^{z}$ in (\ref{tildeCBdef3}) may be replaced by the 
factored equation
\begin{eqnarray} \label{GGz1}
\tilde{\CB P}^{z}
&=&-i {\bmi I}_L \otimes\left[\left({\bmi I}_6\otimes 
{\bf DFT}({\bf k}_{x}) \right){\bmi  D}_{z}({\bf k}_{x},{k}_{y})\right]
 {\CB Q}^{z}\nonumber\\
 &&\qquad\qquad\qquad
 \times
 {\bmi I}_L \otimes \left[
{\bmi  D}^{T}_{z}(-{\bf k}_{x},-{k}_{y})\left({\bmi I}_6\otimes
{\bf DFT}^{T}(-{\bf k}_{x})  \right)\right]
\end{eqnarray} 
where
\begin{equation} \label{CBQdef}
{\CB Q}^{z}=\left[{\begin{array}{cccc}
{\bf 0}_{6N_k}&{\bmi Q}^{z}_{1 
\leftarrow   2}&\cdots &{\bmi Q}^{z}_{1 \leftarrow   L}\\ {\bmi Q}^{z}_{2 
\leftarrow   1}&{\bf 0}_{6N_k}&&\vdots\\ 
\vdots&&\ddots&{\bmi Q}^{z}_{L-1 \leftarrow   L}\\ {\bmi Q}^{z}_{L \leftarrow   
1}&\cdots &{\bmi Q}^{z}_{L \leftarrow   L-1}&{\bf 0}_{6N_k}\end{array}}\right].
\end{equation}
We employ (\ref{Qdefg}) for the partitions ${\bmi Q}^{z}_{i \leftarrow   
j}$ in (\ref{CBQdef}) to get
\begin{equation} \label{Qzcomp}
{\CB Q}^{z}=\left[
{\begin{array}{ccccc}{\begin{array}{cc}{\bf 0}&{\bf 0}\\
{\bf 0}&{\bf 0}\end{array}}&{\begin{array}{cc}{\bf 0}&{\bf 0}\\
 {\bf E}& {\bf 0}\end{array}}&{\begin{array}{cc}{\bf 0}&{\bf 0}\\
{ {\bf E}}^{2}&{\bf 0}\end{array}}&\cdots &{\begin{array}{cc}{\bf 0}&{\bf 0}\\
{ {\bf E}}^{L-1}&{\bf 0}\end{array}}\\
{\begin{array}{cc}{\bf 0}& {\bf E}\\
 {\bf 0}&{\bf 0}\end{array}}&{\begin{array}{cc}{\bf 0}&{\bf 0}\\
{\bf 0}&{\bf 0}\end{array}}&{\begin{array}{cc}{\bf 0}&{\bf 0}\\
 {\bf E}& {\bf 0}\end{array}}&\cdots &{\begin{array}{cc}{\bf 0}&{\bf 0}\\
{ {\bf E}}^{L-2}&{\bf 0}\end{array}}\\
{\begin{array}{cc}{\bf 0}&{ {\bf E}}^{ 2}\\
{\bf 0}&{\bf 0}\end{array}}&&\ddots&& \vdots\\
 \vdots&&&{\begin{array}{cc}{\bf 0}&{\bf 0}\\
{\bf 0}&{\bf 0}\end{array}}&{\begin{array}{cc}{\bf 0}&{\bf 0}\\
 {\bf E}& {\bf 0}\end{array}}\\
{\begin{array}{cc}{\bf 0}&{ {\bf E}}^{ L-1}\\
{\bf 0}&{\bf 0}\end{array}}&{\begin{array}{cc}{\bf 0}&{ {\bf E}}^{ L-2}\\
{\bf 0}&{\bf 0}\end{array}}& \cdots &{\begin{array}{cc} {\bf 0}&{ {\bf E}}^{}\\
{\bf 0}&{\bf 0}\end{array}}&{\begin{array}{cc} {\bf 0}&{\bf 0}\\
{\bf 0}&{\bf 0}\end{array}}\end{array}}
\right]
\end{equation}
We can reduce the floating point operation count of $\tilde{\CB 
P}{}^{z}\cdot \tilde{\CB F}$, with $\tilde{\CB 
P}{}^{z}$ given by (\ref{GGz1}), by factorizing  the ${\CB Q}^{z}$ factor, 
defined by (\ref{Qzcomp}).

We demonstrate the plane wave propagator factorization technique, by 
factoring ${\CB Q}^{z}$ when the lattice has $3$ rows, i.e. $L=3$, so that (\ref{Qzcomp}) reduces to
\begin{equation} \label{Qz3}
{\CB Q}^{z}=
\left[{\begin{array}{cccccc} 
{\bf 0}&{\bf 0}&{\bf 0}&{\bf 0}&{\bf 0}&{\bf 0}\\
{\bf 0}&{\bf 0}& {\bf E}& {\bf 0}&{\bf E}^{2}&{\bf 0}\\
{\bf 0}& {\bf E}& {\bf 0}&{\bf 0}&{\bf 0}&{\bf 0}\\
{\bf 0}&{\bf 0}&{\bf 0}&{\bf 0}& {\bf E}& {\bf 0}\\
{\bf 0}&{\bf E}^{2}&{\bf 0}& {\bf E}& {\bf 0}&{\bf 0}\\
{\bf 0}&{\bf 0}&{\bf 0}&{\bf 0}&{\bf 0}&{\bf 0}\end{array}}\right]
\end{equation}
which may then be factored as
\begin{equation} \label{Qz3fac}
{\CB Q}^{z}={\left[{\begin{array}{cccccc} 
{\bmi I}&{\bf 0}&{\bf 0}&{\bf 0}&{\bf 0}&{\bf 0}\\
{\bf 0}&{\bmi I}&{\bf 0}&{\bf 0}&{\bf 0}&{\bf 0}\\
-\, {\bf E}& {\bf 0}&{\bmi I}&{\bf 0}&{\bf 0}&{\bf 0}\\
{\bf 0}&{\bf 0}&{\bf 0}&{\bmi I}&{\bf 0}&{\bf 0}\\
{\bf 0}&{\bf 0}&-\, {\bf E}& {\bf 0}&{\bmi I}&{\bf 0}\\
{\bf 0}&{\bf 0}&{\bf 0}&{\bf 0}&{\bf 0}&{\bmi I}\end{array}}\right]}^{-1}
{\left[{\begin{array}{cccccc} {\bmi I}&{\bf 0}&{\bf 0}&{\bf 0}&{\bf 0}&{\bf 0}\\
{\bf 0}&{\bmi I}&{\bf 0}&-\, {\bf E}& {\bf 0}&{\bf 0}\\
{\bf 0}&{\bf 0}&{\bmi I}&{\bf 0}&{\bf 0}&{\bf 0}\\
{\bf 0}&{\bf 0}&{\bf 0}&{\bmi I}&{\bf 0}&-\, {\bf E}\\
 {\bf 0}&{\bf 0}&{\bf 0}&{\bf 0}&{\bmi I}&{\bf 0}\\
{\bf 0}&{\bf 0}&{\bf 0}&{\bf 0}&{\bf 0}&{\bmi I}\end{array}}\right]}^{-1}
\left[{\begin{array}{cccccc} {\bf 0}&{\bf 0}&{\bf 0}&{\bf 0}&{\bf 0}&{\bf 0}\\
{\bf 0}&{\bf 0}& {\bf E}& {\bf 0}&{\bf 0}&{\bf 0}\\
{\bf 0}& {\bf E}& {\bf 0}&{\bf 0}&{\bf 0}&{\bf 0}\\
{\bf 0}&{\bf 0}&{\bf 0}&{\bf 0}& {\bf E}& {\bf 0}\\
{\bf 0}&{\bf 0}&{\bf 0}& {\bf E}& {\bf 0}&{\bf 0}\\
{\bf 0}&{\bf 0}&{\bf 0}&{\bf 0}&{\bf 0}&{\bf 0}\end{array}}\right]
\end{equation}

Let
\begin{equation} \label{zsource}
\left[\begin{array}{c}
{\bmi \Sigma}^{\rho^{\prime}}_{\uparrow}\\
{\bmi \Sigma}^{\rho^{\prime}}_{\downarrow}
\end{array}  \right]\equiv
-i\,
{\bmi  D}^{T}_{z}(-{\bf k}_{x},-{k}_{y})\left({\bmi I}_6\otimes
{\bf DFT}^{T}(-{\bf k}_{x})  \right)\cdot \tilde{\CB F}_{\rho^{\prime}}
\end{equation}
denote the $\left[{\stackrel{\mbox{\small upgoing}}{{}_{\mbox{\small 
downgoing}}}}\right]$ planewave radiation from $\mbox{row}_{\rho^{\prime}}$
for each  $1 \le {\rho^{\prime}} \le 3$.

 
Then the action 
\[
{\CB Q}^{z}
 {\bmi I}_L \otimes \left[
{\bmi  D}^{T}_{z}(-{\bf k}_{x},-{k}_{y})\left({\bmi I}_6\otimes
{\bf DFT}^{T}(-{\bf k}_{x})  \right)\right]\cdot \tilde{\CB F}
\]
may be computed using the 
factored representation of (\ref{Qz3fac}), computing the inverses using 
forward and backward substitution. The result of the calculation may be 
summarized as:
\begin{equation} \label{zmarch}
\left.\begin{array}{cc}
{\mit \Delta} { \bmi V}_{\uparrow }^{  {\rho^{\prime}}}  = {\bf E}\left({{ \bmi
\Sigma}_{ \uparrow }^{{\rho^{\prime}}+1}+ {\mit \Delta} {\bmi  V}_{\uparrow }^{  
{\rho^{\prime}}+1}}\right)&{\rho^{\prime}}=1,2\\
{\mit \Delta} {\bmi  V}_{\downarrow }^{  {\rho^{\prime}}} ={\bf E}\left({{ \bmi
\Sigma}_{ \downarrow }^{{\rho^{\prime}}-1}+ {\mit \Delta} { \bmi V}_{\downarrow }^{{\rho^{\prime}}-1}}
\right)& {\rho^{\prime}}=2,3
\end{array}\right\}
\end{equation} 
where ${\mit \Delta} { \bmi V}_{\uparrow }^{  {\rho^{\prime}}}$(${\mit \Delta} { \bmi 
V}_{\downarrow }^{  {\rho^{\prime}}}$) is the upward(downward) incident field on 
$\mbox{row}_{\rho^{\prime}}$ of the lattice. It is not too difficult to generalize 
(\ref{zmarch}) to a lattice with $L$ rows, in which case we get
\begin{equation} \label{zmarchL}
\left.\begin{array}{cl}
{\mit \Delta} { \bmi V}_{\uparrow }^{  {\rho^{\prime}}}  = {\bf E}\left({{ \bmi
\Sigma}_{ \uparrow }^{{\rho^{\prime}}+1}+ {\mit \Delta} {\bmi  V}_{\uparrow }^{  
{\rho^{\prime}}+1}}\right)&{\rho^{\prime}}=1,L-1\\
{\mit \Delta} {\bmi  V}_{\downarrow }^{  {\rho^{\prime}}} ={\bf E}\left({{ \bmi
\Sigma}_{ \downarrow }^{{\rho^{\prime}}-1}+ {\mit \Delta} { \bmi V}_{\downarrow }^{{\rho^{\prime}}-1}}
\right)& {\rho^{\prime}}=2,L
\end{array}\right\}
\end{equation} 

\begin{samepage}
Given that
\[
\left[\begin{array}{c}
{\mit \Delta} {\bmi  V}_{\downarrow }^{  \rho}\\
{\mit \Delta} {\bmi  V}_{\uparrow }^{  \rho}
\end{array}  \right]
\]
\end{samepage}
denotes the $\left[{\stackrel{\mbox{\small downgoing}}{{}_{\mbox{\small 
upgoing}}}}\right]$
incident plane wave radiation on $\mbox{row}_{\rho}$ of the lattice, we 
synthesize the  $6 N_k$ displacement/traction vector using
\begin{equation} \label{zsynthesis}
\left[\begin{array}{c}
\tilde{\bmi u}_{  \rho}\\
\tilde{\bmi t}_{  \rho}
\end{array}  \right]
=
\left({\bmi I}_6\otimes 
{\bf DFT}({\bf k}_{x}) \right){\bmi  D}_{z}({\bf k}_{x},{k}_{y})
\left[\begin{array}{c}
{\mit \Delta} {\bmi  V}_{\downarrow }^{  \rho}\\
{\mit \Delta} {\bmi  V}_{\uparrow }^{  \rho}
\end{array}\right]
\end{equation}
operation. Finally, we downsample or decimate this result, to get
\[
\left[\begin{array}{c}
{\bmi u}_{  \rho}\\
{\bmi t}_{  \rho}
\end{array}  \right]=
{\bmi \Phi}^{-1}\left[\begin{array}{c}
\tilde{\bmi u}_{  \rho}\\
\tilde{\bmi t}_{  \rho}
\end{array}  \right],
\]
the $6$ component displacement/traction vector  at each of the $N_x$ 
lattice points in $\mbox{row}_{\rho}$ of the lattice.

In summary, if we implement the $6 N_x L$ matrix/vector mapping
${\CB P}{}^{z}\cdot {\CB F}$ of (\ref{CBdef03}) using 
(\ref{zsource})--(\ref{zsynthesis}), the floating point operation count 
is asymptotically
\begin{equation} \label{N5def}
N_5=2\,(L-1)\,\left[ 39 N_k + 6 N_k \log_2 N_k \right]
\end{equation}
This algorithm is an elastic wave generalization of the acoustic wave $(\omega, k)$ {\em phase shift 
method\/} of {\em Gazdag\/}~(1981). However, our algorithm has been set 
up with the intention of extending it to a plane stratified background 
model.

%%%%%%%%%%%%%%%%%%%%%%%%%%%%%%%%%%%%%%%%%%%%%%%%%%%%%%%%%%%%%%%%%%%%%%%%%%%%%%%%%%%%%%%%%%
\subsection{ Implementing $\mbox{\CBM  P}\,{}^{x}\cdot\mbox{\CBM  F}$}

We now find a representation of the $6 N_x \times 6 N_x$ partition  ${\CB P}_{\rho \leftarrow 
\rho}$ of array ${\CB P}^{x}$ in (\ref{CBdef4}) using (\ref{superPx}) for the 
$6 \times 6$ component partitions. 

From (\ref{CBdef2}) with $\rho=\rho^{\prime}$, 
and using (\ref{Gijdef}) for the definition of the  partitions we have
\begin{eqnarray} \label{CBdef5}
{\CB P}_{\rho \leftarrow   \rho}=
\left[{\begin{array}{cccc}
{\bf 0}_{6}&\rf{\bmi P}\!{}^{x}({x_1},{x_2})&\cdots 
&\rf{\bmi P}\!{}^{x}({x_1},{x_{N_x}})\\
\rf{\bmi P}\!{}^{x}({x_2},{x_1})&{\bf 0}_{6}&&\vdots\\
\vdots&&\ddots&\rf{\bmi P}\!{}^{x}(x_{N_x-1} , {x_{N_x}})\\
\rf{\bmi P}\!{}^{x}(x_{N_x},x_1)&\cdots &
\rf{\bmi P}\!{}^{x}(x_{N_x},x_{N_x-1})&{\bf 0}_{6}
\end{array}}\right]
\end{eqnarray}


By employing the propagator of (\ref{superPx}) for the partitions of  (\ref{CBdef5}),
 we get
\begin{samepage}
\begin{eqnarray}	\label{GGx0}
\lefteqn{\rf{\CB P}_{\rho \leftarrow   \rho} \equiv}\\
&&
 -i{\bmi I}_{N_x} \otimes
\left[\left({\bmi I}_6\otimes \overline{\bmi F}(0)  \right)
 {\bmi  D}_{x}({\bf k}_{z},{k}_{y})\right]
 {\CB Q}^{x}_{\rho \leftarrow \rho}
{\bmi I}_{N_x} \otimes \left[{\bmi  D}_{x}^{T}(-{\bf k}_{z},-{k}_{y})
\,\left({\bmi I}_6\otimes \overline{\bmi F}^{T}(0)  \right)
\right]\nonumber
\end{eqnarray}
\end{samepage}
where:
\begin{itemize}
\begin{samepage}
\item the  $(6 N_x N_k) \times (6 N_x N_k)$, {\em transverse\/} 
plane wave propagator ${\CB Q}^{x}_{\rho \leftarrow \rho}$ is  defined as
\begin{equation} \label{CBQxdef}
{\CB Q}^{x}_{\rho \leftarrow   \rho}=
\left[{\begin{array}{cccc}
{\bf 0}_{6N_k}&{\bmi Q}^{x}({x}_1,{x}_2)&\cdots 
&{\bmi Q}^{x}({x}_1,{x}_{N_x})\\
{\bmi Q}^{x}({x}_2,{x}_1)&{\bf 
0}_{6N_k}&&\vdots\\
\vdots&&\ddots&{\bmi Q}^{x}({x}_{N_x -1},{x}_{N_x})\\
{\bmi Q}^{x}({x}_{N_x},{x}_{1})&\cdots &
{\bmi Q}^{x}({x}_{N_x},{x}_{N_x-1})&{\bf 0}_{6N_k}
\end{array}}\right]
\end{equation}
\end{samepage}
where the  $6 N_k \times 6 N_k$ plane wave propagator partition, 
${\bmi Q}^{x}({x}_{\chi},{x}_{\chi^{\prime}}), \ 1 \le \chi, \chi^{\prime} \le N_x$ in 
(\ref{CBQxdef}) 
is given by  
\begin{equation} \label{Qxdefg}
{\bmi Q}^{x}({x}_{\chi},{x}_{\chi^{\prime}})=
\left[{\begin{array}{cc}{\bf 0}_{3N_k}&H(\chi - \chi^{\prime}){\bf E}_{x}^{|\chi - \chi^{\prime}|}\\
H(\chi^{\prime}-\chi){\bf E}_{x}^{|\chi - \chi^{\prime}|}&{\bf 0}_{3N_k}\end{array}}\right]
\end{equation}
with the $3 N_k$-diagonal phase shift array ${\bf E}_{x}$ defined by
 \begin{equation} \label{Exdef}
{\bf E}_{x}=\left[{\begin{array}{ccc}
{\rm e}^{i{\bf k}_{x,P}\,{\mit \Delta}x}&{\bf 0}_{N}&{\bf 0}_{N}\\ 
{\bf 0}_{N}&{\rm e}^{ 
i{\bf k}_{x,S}\,{\mit \Delta}x}&{\bf 0}_{N}\\ {\bf 0}_{N}&{\bf 0}_{N}&{\rm e}^{ 
i{\bf k}_{x,H}\,{\mit \Delta}x}\end{array}}\right] 
\end{equation}
and  ${\bf k}_{x,c}, \ c\in\{ P,S,H\}$ was given by (\ref{gbarKx}), so that
\begin{equation} \label{Qxcomp}
{\CB Q}^{x}_{\rho \leftarrow   \rho}=\left[
{\begin{array}{ccccc}{\begin{array}{cc}{\bf 0}&{\bf 0}\\
{\bf 0}&{\bf 0}\end{array}}&{\begin{array}{cc}{\bf 0}&{\bf 0}\\
 {\bf E}_{x}& {\bf 0}\end{array}}&{\begin{array}{cc}{\bf 0}&{\bf 0}\\
{\bf E}_{x}^{2}&{\bf 0}\end{array}}&\cdots &{\begin{array}{cc}{\bf 0}&{\bf 0}\\
{\bf E}_{x}^{N-1}&{\bf 0}\end{array}}\\
{\begin{array}{cc}{\bf 0}& {\bf E}_{x}\\
 {\bf 0}&{\bf 0}\end{array}}&{\begin{array}{cc}{\bf 0}&{\bf 0}\\
{\bf 0}&{\bf 0}\end{array}}&{\begin{array}{cc}{\bf 0}&{\bf 0}\\
 {\bf E}_{x}& {\bf 0}\end{array}}&\cdots &{\begin{array}{cc}{\bf 0}&{\bf 0}\\
{\bf E}_{x}^{N-2}&{\bf 0}\end{array}}\\
{\begin{array}{cc}{\bf 0}&{\bf E}_{x}^{ 2}\\
{\bf 0}&{\bf 0}\end{array}}&&\ddots&& \vdots\\
 \vdots&&&{\begin{array}{cc}{\bf 0}&{\bf 0}\\
{\bf 0}&{\bf 0}\end{array}}&{\begin{array}{cc}{\bf 0}&{\bf 0}\\
 {\bf E}_{x}& {\bf 0}\end{array}}\\
{\begin{array}{cc}{\bf 0}&{\bf E}_{x}^{ N-1}\\
{\bf 0}&{\bf 0}\end{array}}&{\begin{array}{cc}{\bf 0}&{\bf E}_{x}^{ N-2}\\
{\bf 0}&{\bf 0}\end{array}}& \cdots &{\begin{array}{cc} {\bf 0}&{\bf E}_{x}\\
{\bf 0}&{\bf 0}\end{array}}&{\begin{array}{cc} {\bf 0}&{\bf 0}\\
{\bf 0}&{\bf 0}\end{array}}\end{array}}
\right]
\end{equation}


\item the $6 N_k \times 6 N_k$ array ${\bmi  D}_{x}({\bf k}_{z},{k}_{y})$ was defined in (\ref{Dxsuperdef})

\item the $6 N_k$-vector $\overline{\bmi F}(0)$ was defined in (\ref{Fbardef})
\end{itemize}

With $\rf{\CB P}_{\rho \leftarrow   \rho}$ given by (\ref{GGx0}), and  ${\CB  F}_{\rho}$ given by
(\ref{Frho}), we can reduce the cost of the $6 N_x$ matrix/vector 
products
\[
\rf{\CB P}_{\rho \leftarrow   \rho}\cdot{\CB  F}_{\rho}, \qquad 1 \le \rho \le L
\]
by factoring the plane wave propagator ${\CB Q}^{x}_{\rho \leftarrow   \rho}$
component of $\rf{\CB P}_{\rho \leftarrow   \rho}$ that was defined in 
(\ref{Qxcomp}).

Let us denote the left(right) going plane waves from a lattice 
point in $\mbox{column}_\chi$ of $\mbox{row}_\rho$ by ${\bmi 
\Sigma}^{(\rho,\chi)}_{\leftarrow}$(${\bmi \Sigma}^{(\rho,\chi)}_{\rightarrow}$) for 
$1 \le \rho \le L$, $1 \le \chi \le N_x$, then
\begin{eqnarray} \label{LRsource}
\left[\begin{array}{c}
{\bmi \Sigma}^{(\rho,\chi)}_{\leftarrow}\\
{\bmi \Sigma}^{(\rho,\chi)}_{\rightarrow}
\end{array}\right]&=&
{\bmi  D}_{x}^{T}(-{\bf k}_{z},-{k}_{y})
\,\left({\bmi I}_6\otimes \overline{\bmi F}^{T}(0)  \right){\bmi J}_6 
\bar{\bmi f}_{N_x(\rho-1)+\chi}\nonumber\\
&=&
{\bmi  D}_{x}^{T}(-{\bf k}_{z},-{k}_{y})
\,\left({\bmi I}_6\otimes \overline{\bmi F}^{T}(0)  \right)
{\bmi f}_{N_x(\rho-1)+\chi}
\end{eqnarray}
 By analogy with (\ref{zmarchL}) we have
\begin{equation} \label{xmarchL}
\begin{array}{ll}
\left.\begin{array}{cl}
{\mit \Delta} { \bmi V}_{\leftarrow }^{(\rho,\chi)}  = {\bf E}_{x}\left({{ \bmi
\Sigma}_{ \leftarrow }^{(\rho,\chi+1)}+ {\mit \Delta} {\bmi  V}_{\leftarrow }^{  
(\rho,\chi+1)}}\right)&\chi=1,N_x-1\\
{\mit \Delta} {\bmi  V}_{\rightarrow }^{(\rho,\chi)} ={\bf E}_{x}\left({{ \bmi
\Sigma}_{ \rightarrow }^{(\rho,\chi-1)}+ {\mit \Delta} { \bmi V}_{\rightarrow 
}^{(\rho,\chi-1)}}
\right)& \chi=2,N_x
\end{array}\right\}&\rho=1,L
\end{array}
\end{equation} 
where  ${\mit \Delta} { \bmi V}_{\leftarrow }^{(\rho,\chi)}$(${\mit \Delta} { \bmi 
V}_{\rightarrow }^{(\rho,\chi)}$) is the leftward(rightward) incident field on 
$(\mbox{row}_\rho,\mbox{column}_\chi)$ of the lattice.
Finally, we synthesize the components of displacement and traction at
$(\mbox{row}_\rho,\mbox{column}_\chi)$ from the incident plane waves, using
\begin{eqnarray} \label{utxsyn}
\left[\begin{array}{c}
{\bmi u}^{(\rho,\chi)}\\
{\bmi t}^{(\rho,\chi)}
\end{array}\right]=
\left({\bmi I}_6\otimes \overline{\bmi F}(0)  \right)
 {\bmi  D}_{x}({\bf k}_{z},{k}_{y})
 \left[\begin{array}{c}
{\mit \Delta}{\bmi V}^{(\rho,\chi)}_{\rightarrow}\\
{\mit \Delta}{\bmi V}^{(\rho,\chi)}_{\leftarrow}
\end{array}\right], \quad \rho=1,L, \quad \chi=1,N_x
\end{eqnarray}


%%%%%%%%%%%%%%%%%%%%%%%%%%%%%%%%%%%%%%%%%%%%%%%%%%%%%%%%%%%%%%%%%%%%%%%%%%%%%%%%%%%%%%%%%
The computations (\ref{LRsource})--(\ref{utxsyn}) may be summarized by the matrix formula
\[
\left[ \begin{array}{c}{\bmi u}\\ {\bmi t}\end{array} \right]=
\rf{\CB P}\!{}^{x}\cdot {\CB  F}
\]
with
\begin{samepage}
\begin{eqnarray}	\label{GGx1}
\lefteqn{\rf{\CB P}\!{}^{x} \equiv -{i}{\bmi I}_L \otimes\left( \right.}\\
&&
\left.
{\bmi I}_{N_x} \otimes
\left[\left({\bmi I}_6\otimes \overline{\bmi F}(0)  \right)
 {\bmi  D}_{x}({\bf k}_{z},{k}_{y})\right]
 {\CB Q}^{x}_{\rho \leftarrow \rho}
{\bmi I}_{N_x} \otimes \left[{\bmi  D}_{x}^{T}(-{\bf k}_{z},-{k}_{y})
\,\left({\bmi I}_6\otimes \overline{\bmi F}^{T}(0)  \right)
\right] \right)\nonumber
\end{eqnarray}
\end{samepage}
an explicit representation of the array ${\CB P}^{x}$ introduced in (\ref{CBdef4}),
that follows after using (\ref{GGx0}) for the diagonal blocks ${\CB P}^{x}_{\rho \leftarrow 
\rho}, \ 1 \le \rho \le L$.
The asmyptotic floating point operation count of eqns. (\ref{LRsource})--(\ref{utxsyn})
is
\begin{equation} \label{N6def}
N_6=78\,L\,(N_x-1)  \overline{N}_k
\end{equation} 



%%%%%%%%%%%%%%%%%%%%%%%%%%%%%%%%%%%%%%%%%%%%%%%%%%%%%%%%%%%%%%%%%%%%%%%%%%%%%%%%%%%%%%%%%
%%%%%%%%%%%%%%%%%%%%%%%%%%%%%%%%%%%%%%%%%%%%%%%%%%%%%%%%%%%%%%%%%%%%%%%%%%%%%%%%%%%%%%%%%
%%%%%%%%%%%%%%%%%%%%%%%%%%%%%%%%%%%%%%%%%%%%%%%%%%%%%%%%%%%%%%%%%%%%%%%%%%%%%%%%%%%%%%%%%
\subsection{ Implementing $\mbox{\CBM  P}\cdot\mbox{\CBM  F}$}


Finally, in the culmination of the chapter, we use (\ref{GGz1}) and (\ref{GGx1}) in 
(\ref{CBGsplit}) to get the representation of 
the {\em lattice propagator\/} $\rf{\CB P}$, which may be written as
\[
\!\!\!\!\!\!\fbox{
\begin{minipage} {146.0mm}
\begin{samepage}
\begin{eqnarray}\label{GGlattice}
\!\!\!\!\!\!\!\!\!
\rf{\CB P} &\equiv&
\rf{\CB P}\!{}^{z}\left[ \cdot \right]+ \rf{\CB  P}\!{}^{x}\nonumber\\
&&\!\!\!\!\!\!= -i{\bmi \Phi}{}^{-1}\left[  {\bmi I}_L \otimes\left[\left({\bmi I}_{6}\otimes 
{\bf DFT}({\bf k}_{x}) \right){\bmi  D}_{z}({\bf k}_{x},{k}_{y})\right]
 {\CB Q}^{z}\right.\nonumber\\
&&\qquad\qquad\qquad\qquad\qquad
\times \left.
 {\bmi I}_L \otimes \left[
{\bmi  D}^{T}_{z}(-{\bf k}_{x},-{k}_{y})\left({\bmi I}_{6}\otimes
{\bf DFT}^{T}(-{\bf k}_{x})  \right)  {\bmi \Phi}\left[ \cdot \right] \right]\right]\nonumber\\
&&
 -{i}{\bmi I}_L \otimes\left\{ {\bmi I}_{N_x} \otimes
\left[\left({\bmi I}_{6}\otimes \overline{\bmi F}(0)  \right)
 {\bmi  D}_{x}({\bf k}_{z},{k}_{y})\right]
 {\CB Q}{}^{x}_{\rho \leftarrow \rho}
\right.\nonumber\\
&&\qquad\qquad\qquad\qquad\qquad
\left.
\times{\bmi I}_{N_x} \otimes \left[{\bmi  D}_{x}^{T}(-{\bf k}_{z},-{k}_{y})
\,\left({\bmi I}_{6}\otimes \overline{\bmi F}^{T}(0)  \right)
\right] \right\}
\end{eqnarray}
\end{samepage}
\vspace{1 mm}
\end{minipage}
}
\]
Note that we have incorporated to upsampling and downsampling operators 
within (\ref{GGlattice}), so that $\rf{\CB P}$ operates on $6 L 
N_x$-vectors.


In (\ref{xmarchL}), we have a left/right sweep in each of the $L$ rows of 
the lattice, where as (\ref{zmarchL}) is composed of a single up/down 
sweep. An accurate measure of the number of floating point operations 
embodied in  (\ref{zmarchL})[(\ref{xmarchL})] may be ascertained by counting 
the number of applications of the vertical[horizontal] phase shift arrays 
${\bf E}$[${\bf E}_{x}$]. For our $L-$row, $N_x-$column lattice, we have
$2(L-1)$ applications of ${\bf E}$ and $2 L (N_x-1)$ applications of 
${\bf E}_{x}$, where ${\bf E}$ is a $3 N_k$-element diagonal array and 
${\bf E}_{x}$ is a $3 \overline{N}$-element diagonal array, with 
$\overline{N}_k\le 
{N}_k$. 

We generally use less plane waves in the transverse phase shifts of ${\bf E}_{x}$
than the vertical phase shifts of ${\bf E}$ because all points in a 
given row have the same depth coordinate $z$ and consequently, there is no possibility of 
{\em aliasing\/}
caused by 
undersampling of the complex exponential $\exp (i k_z {\mit \Delta}z)$ 
because ${\mit \Delta}z\equiv 0$ when phase shifting between points 
within the same row. However, when we compute the vertical phase shifts, 
we must ensure that we have employed an adequate $k_x-$wavenumber 
bandwidth $2 {\mit \Omega}$ and sampling rate ${\mit \Delta}k_x$ so that the DFT arrays 
${\bf DFT}$ in (\ref{GGlattice}) correctly model the transverse phase shifts between 
lattice points at {\em opposite ends  of  adjacent rows in the 
lattice\/}. 
At the Nyquist sampling rate, the number of lattice points within a given 
row $N_x$ is equal to the number of plane waves $N_k$ used to phase shift between 
rows.
 However, in numerical applications of the techniques described 
here, we have found it necessary to employ a considerable amount of 
{\em oversampling\/}, so that
\[
N_k = {N}_{\mit \Delta } {N}_{L } N_x.
\]
which was defined above in (\ref{Nkdef}). In the numerical computations 
of chapter \ref{variational}, ${N}_{\mit \Delta }=8$, ${N}_{L }=4$, so that
\[
N_k = 32  N_x
\]
i.e., $32$ times the number of plane 
waves that would be necessary from Shannon's sampling theorem (Jerri, 1977; Brigham, 1988). 
We compute 
the relevant quantities at the actual grid points by {\em decimating\/} 
or {\em downsampling\/} 
the results at the fictitious oversampled lattice points (Alg. 3.1).
Usually, we have taken $\overline{N}_k=N_k/2$ or $\overline{N_k}=N_k/4$, where 
$N_k=2048$ or $N_k=4096$ is the number of oversampled plane waves 
employed in the vertical phase shifts between rows and $\overline{N}_k$ is the 
corresponding number of plane waves in the transverse shifts between 
lattice points within the same row.

The advantages of the representation (\ref{GGlattice}) are twofold:
\begin{enumerate}
\item the action of (\ref{GGlattice}) on a secondary source field 
${\CB F}$ with $6$-vector components ${\bmi f}_j$ at each of the $N_x\times L$ lattice 
points of Fig\ref{lattice},  may be computed in the order of 
\[
N_2\approx2\,(L-1)\,\left[ 39 N_k + 6 N_k \log_2 N_k \right]+ 
78\,L\,(N_x-1)  \overline{N}_k
\]
floating point operations, which is often  far less arithmetic than 
the ostensible 
\[
N_1\approx36 N_x^2 \times L^2 - 6 N_x \times L
\]
 floating point operation count. However,  the oversampling requirement
 (\ref{Nkdef}),
 erodes the numerical advantage of $N_2$ over $N_1$.  
\item  (\ref{GGlattice}) is readily extensible to 
the important case where the lattice is superimposed on a plane 
stratified half space. We consider this problem in chapter \ref{stratG}.
\end{enumerate} 


In summary,
the key steps in achieving the goal expressed by item-({\bf 1}) above is to:
\begin{enumerate} 
\item factor the plane wave propagators ${\CB Q}^{z}$ and ${\CB Q}{}^{x}$ 
in (\ref{GGlattice}) into the product of {\em sparse\/} matrices
\item employ the {\em Fast Fourier Transform\/}(FFT) algorithm to compute the
action of the {\em Discrete Fourier Transform\/}(DFT) matrix ${\bf DFT}$ 
from (\ref{DFTdef}) on a field vector. The FFT algorithm factors ${\bf DFT}$ into 
a product of sparse arrays so that the cost of applying ${\bf DFT}$ to an 
$N_k-$vector is reduced from $O(N_k^2)$ to $O(N_k \log_2 N_k)$ floating point 
operations. 
\end{enumerate}
 

In closing this chapter, we note that although we have assumed that the 
secondary source vector ${\CB F}$ in (\ref{GFdef0}), has $6-$vector components
${\bmi f}_j, \ 1 \le j \le L N_x$, it will be shown in chapter \ref{Trep}, that
${\bmi f}_j, \ 1 \le j \le L N_x$ is most conveniently expressed as an 
$8-$vector. The reason we have employed $6-$vectors here is related to 
the fact that the partial wave displacement/traction states  used to 
construct the Green's tensor have $6$ components, i.e. $3 $ components 
of displacement $(u_z,u_x,u_y)$, and $3$ components of traction
$(\tau_{zz},\tau_{xz},\tau_{yz})$. In chapter \ref{Trep}, 
we employ  integration by parts, so that the partial wave components of 
the Green's tensor that propagates the field between successive scattering 
interactions, are augmented by the $(\partial_{x}u_x,\partial_{x}u_y)$ 
components. Consequently, the {\em effective\/} 
secondary source vector $\hat{\CB F}$ has $8-$vector components, say,
$\hat{\bmi f}_j, \ 1 \le j \le L N_x$. Some of the results presented here 
will then require minor modification, and these modifications will be 
outlined in chapter \ref{Trep}.




